\documentclass[11pt]{ctexart}
\usepackage[a4paper,margin=1in]{geometry}
\usepackage{graphicx}
\usepackage{xcolor}
\usepackage{hyperref}
\definecolor{refgray}{gray}{0.45}
\title{\textbf{Market Views｜全球投行叙事汇编}\\[0.4em]{\large AI投资挤压传统软件支出，中国科技配置创历史新高，全球市场呈现结构性分化与政策预期分歧}}
\author{KC桌面 自动生成}
\date{260725}
\begin{document}
\maketitle
\tableofcontents
\newpage
\section{一页摘要}
\begin{itemize}
  \item 中国产能结构分化加剧：新兴领域（半导体、新能源车出口、储能）增长强劲，传统部门（水泥、光伏内需）承压，外部需求成为消化国内产能的关键渠道。
  \item 欧洲央行加息本身不会导致欧元走弱，但欧元区内部财政分化（法意利差走阔）是需警惕的尾部风险，当前利率水平尚未触及货币转弱的阈值。
  \item AI基础设施投资持续超预期，WFE市场2027年有望突破2000亿美元，但设备商受益程度因DRAM敞口、英特尔资本开支结构变化而显著分化。
  \item 企业资本支出正从传统软件采购转向AI基础设施、存储和内存，导致IBM等传统软件厂商短期业绩承压，但推迟的交易属于需求延迟而非消失。
  \item 全球数据中心容量到2030年将翻倍以上，美国占净增量的60-70\%，但供给端加速与需求可见度之间的差距正在制造新的分化。
  \item 中国公募基金在2026年第二季度大幅增配科技板块，电子、通信持仓占比和超配比例均创历史新高，科技主题已成为核心布局方向。
  \item GLP-1竞争从减重疗效转向代谢获益与便利性，口服剂型（小分子vs肽类）的依从性优势存在分歧；手术机器人长期渗透空间广阔，短期担忧提供窗口。
  \item 全球能源与材料消费遵循“加法”而非“替代”模式，清洁能源占比上升并未导致化石燃料绝对消费下降，可再生能源扩张的核心驱动在于土地与现有装机基础。
\end{itemize}
\section{全球投行叙事汇编}
今日纳入 91 篇投行/券商研究，按 16 个市场、资产与产业主题系统整理。
\newpage
\subsection{宏观与利率：中国产能结构分化与锂价回调的宏观映射}
\textbf{中国6月数据强化了新兴领域（半导体、新能源车出口、储能）与传统部门（水泥、光伏内需）之间的结构性分化，同时锂价从高位回调但市场定价已过度悲观，中间价机制维持稳定。}
\subsubsection*{主流共识}
\begin{itemize}
  \item 中国产能结构分化加剧：新兴领域（半导体、新能源车出口、储能）增长强劲，传统部门（水泥、光伏内需）承压。
  \item 外部需求成为消化国内产能的关键渠道，新能源车出口增速远超国内生产增速。
  \item 锂价短期回调但年内仍累计上涨，储能电池取代电动车成为需求增量主引擎。
\end{itemize}
\subsubsection*{分歧与条件差异}
\begin{itemize}
  \item 半导体出口增长驱动力存在分歧：BARC认为单位价值提升主导，而市场部分观点仍视为数量驱动。
  \item 光伏产业链价格分化：多晶硅库存高企但价格微涨，硅片价格持续下跌，反映不同环节供需格局差异。
  \item 锂价未来走势分歧：HSBC预估未来12个月均价19300美元/吨，而权益市场隐含价格仅13500-14000美元/吨。
  \item 人民币中间价定价中政策调节因子的作用力度：NOM模型显示外部波动大时政策空间收窄，但市场对政策干预预期不一。
\end{itemize}
\subsubsection*{机构观点}
\paragraph{BARC} 中国6月生产与出口数据延续新兴领域与传统部门的分化格局，且外部需求吸收国内产出的角色日益重要。半导体出口额同比增速从5月的110.9\%升至121.9\%，但出口量微降0.3\%，表明增长驱动力转向单位价值提升和产品结构优化。新能源车国内生产端疲弱（混动产量同比-34.2\%），但出口端强劲（纯电出口+81.3\%），凸显内需吸收能力变化。家电出口受益于外部需求回暖，出口价格未跌，显示需求拉动而非降价促销。传统行业如水泥仍以量换价（出口量+109\%，金额仅+88\%）。
{\small\color{refgray}数据：6月集成电路出口额同比+121.9\%，出口量-0.3\%\par}
{\small\color{refgray}数据：新能源车混动产量同比-34.2\%，出口量+105.3\%\par}
{\small\color{refgray}数据：家电空调出口量+11.3\%，金额+10.1\%\par}
{\small\color{refgray}数据：水泥出口量+109\%，金额+88\%\par}
{\small\color{refgray}边际变化：半导体出口从数量驱动转向单位价值提升；新能源车生产与出口分化幅度进一步扩大。\par}
\paragraph{Citi} 中国光伏市场1H26新增装机72.1GW，同比-65.9\%，主要受高基数影响而非需求根本性调整。产业链价格分化：多晶硅价格微涨0.1\%至31.6元/公斤，但库存高达31.8万吨，7月产量10万吨超需求9.1万吨；硅片价格下跌2.3\%至0.84元/W，产能整合启动；组件价格下跌0.4\%至0.71元/W。光伏玻璃行业主动减产，日熔量降至77,750吨（-7.7\%），库存周期47.3天微降。新能效标准2027年生效，有望淘汰低效产能。
{\small\color{refgray}数据：1H26中国光伏新增装机72.1GW，同比-65.9\%\par}
{\small\color{refgray}数据：多晶硅库存31.8万吨，7月产量10万吨超需求9.1万吨\par}
{\small\color{refgray}数据：硅片价格周环比-2.3\%至0.84元/W\par}
{\small\color{refgray}数据：光伏玻璃日熔量较月初-7.7\%至77,750吨\par}
{\small\color{refgray}边际变化：光伏玻璃行业主动减产幅度扩大（-7.7\%），反映盈利压力；新能效标准为中长期供给侧改革信号。\par}
\paragraph{HSBC} 电池级碳酸锂报价自5月高点回落约28\%，但年内仍累计上涨25\%。当前锂矿企业权益定价隐含未来12个月锂均价仅13500-14000美元/吨，较现货低33\%-37\%，而HSBC预估为19300美元/吨，市场定价已过度悲观。需求端，1H26中国电池总销量979GWh（+49\%），储能电池增速83\%远超电动车电池的36\%，AI数据中心成为新需求场景且对价格不敏感。供给端，宁德时代江西矿重启预计2026年贡献4.8万吨LCE，澳洲三矿重启影响2027年。
{\small\color{refgray}数据：锂价自5月高点回落约28\%，年内累计+25\%\par}
{\small\color{refgray}数据：权益市场隐含锂均价13500-14000美元/吨，HSBC预估19300美元/吨\par}
{\small\color{refgray}数据：1H26中国电池总销量979GWh，同比+49\%\par}
{\small\color{refgray}数据：储能电池销量同比+83\%，电动车电池+36\%\par}
{\small\color{refgray}边际变化：储能电池取代电动车成为需求增量主引擎；AI数据中心作为新增需求场景，价格承受力更高。\par}
\paragraph{NOM} 人民币中间价定价框架未现方向性调整，政策调节因子力度温和。NOM模型预测值6.7822，较前一日收盘价高116个基点；计入逆周期因子后为6.7880，较实际中间价低26个基点。模型拆解显示，隔夜美元指数波动幅度超0.3\%时，中间价更多反映市场力量，政策调节空间收窄。模型误差序列为量化政策干预提供了窗口。
{\small\color{refgray}数据：模型预测中间价6.7822，较收盘价高116个基点\par}
{\small\color{refgray}数据：计入逆周期因子后预测6.7880，较实际中间价低26个基点\par}
{\small\color{refgray}数据：隔夜美元指数波动>0.3\%时，模型解释力显著上升\par}
{\small\color{refgray}边际变化：模型预测值较前次下调84个基点，但政策调节因子力度未变，定价框架稳定。\par}
\paragraph{MS} SpaceX 的 Starship 项目在 2023 年至 2026 年中期经历了 13 次飞行测试，其中仅 4 次实现完全达成预期，其余均出现不同程度的助推器或飞船异常。这份MS研报的核心判断是：观察者需要对飞行测试中的异常情况持有现实预期，这些异常可能延迟 SpaceX 在连接和人工智能领域的里程碑，进而影响公司财务表现和报价波动。但报告同时指出，Spa...
{\small\color{refgray}数据：报告以 Starship 飞行 13 的首次尝试为例：四台发动机未能点火，触发自动中止。这一事件导致市场情绪持续调整。然而，分析师认为，一次顺利完成的、仅有轻微异常的测试，可能为报价带来约 5\% 的短期正面波动。在约 115 美元的报价水平上，SpaceX 的航天与 Starlink 业务合计估值约为 136 美元，而当前报价已低于这一分部估值，且未计入任何...\par}
{\small\color{refgray}数据：KC观察： 报告将 SpaceX 的报价拆解为航天、连接、X 与 Grok、企业 AI 四个分部。当前报价仅反映前两个分部的价值，AI 业务尚未被市场定价。这意味着，如果 Starship 测试进展顺利，AI 业务的估值释放可能成为报价上行的额外动力。\par}
\subsubsection*{关键数据}
\begin{itemize}
  \item 6月集成电路出口额同比+121.9\%，出口量-0.3\%
  \item 新能源车混动产量同比-34.2\%，出口量+105.3\%
  \item 1H26中国光伏新增装机72.1GW，同比-65.9\%
  \item 多晶硅库存31.8万吨，7月产量10万吨超需求9.1万吨
  \item 锂价自5月高点回落约28\%，年内累计+25\%
  \item 权益市场隐含锂均价13500-14000美元/吨，HSBC预估19300美元/吨
  \item 1H26中国电池总销量979GWh，同比+49\%，储能电池+83\%
  \item 人民币中间价模型预测6.7822，较收盘价高116个基点
\end{itemize}
\begin{figure}[htbp]
\centering
\includegraphics[width=0.88\linewidth]{figures/fig_109.jpg}
\caption{Figure 1 - Figure 1. PRC weekly polysilicon prices in the week ended Jul 22 © 2026 Citi Inc. No redistribution without Citi's written permission. Figure 2. PRC weekly polysilicon price © 2026 Citi Inc. No redistribution without Citi's written permission. Figure 3. PRC we}
\end{figure}
\begin{figure}[htbp]
\centering
\includegraphics[width=0.88\linewidth]{figures/fig_110.jpg}
\caption{Figure 2 - © 2026 Citi Inc. No redistribution without Citi's written permission. Figure 2. PRC weekly polysilicon price © 2026 Citi Inc. No redistribution without Citi's written permission. Figure 3. PRC weekly wafer price © 2026 Citi Inc. No redistribution without Citi'}
\end{figure}
\begin{figure}[htbp]
\centering
\includegraphics[width=0.88\linewidth]{figures/fig_080.jpg}
\caption{Exhibit 1 - MS does and seeks to do business with companies covered in MS. As a result, investors should be aware that the firm may have a conflict of interest that could affect the objectivity of MS. Investors should consider MS as only a single factor in making their in}
\end{figure}
\begin{figure}[htbp]
\centering
\includegraphics[width=0.88\linewidth]{figures/fig_081.jpg}
\caption{Exhibit 2 - Exhibit 2: Starship Flight Test History √ Achieved ✗ Not achieved / partially achieved — Skipped / not reached Exhibit 3: Collection of Various Falcon 9 Anomalies During Booster Landing Tests \#\# Valuation Methodology and Risks \#\#}
\end{figure}
{\small\color{refgray}本节综合 5 篇研究；涉及 BARC、MS、NOM、Citi、HSBC。}
\newpage
\subsection{欧元利率与汇率：紧缩未必走弱，财政分化才是真正风险}
\textbf{欧洲央行加息本身不会导致欧元走弱，除非紧缩过度拖累增长或引发财政可持续性担忧；当前欧元区利率水平尚未触及历史案例中货币转弱的阈值，但内部利差走阔是需警惕的尾部风险。}
\subsubsection*{主流共识}
\begin{itemize}
  \item 利率上升通常通过利差渠道利好货币，但存在两个例外：紧缩损害增长或引发债务可持续性担忧。
  \item 当前欧元区1年期与5年5年远期利率之差约-78bp，远低于历史案例中货币明显转弱的阈值。
  \item 欧元在2023年7月进入紧缩周期后12个月仅贬值0.7\%，远低于历史平均-0.8\%。
\end{itemize}
\subsubsection*{分歧与条件差异}
\begin{itemize}
  \item 市场对欧央行紧缩程度的定价是否已接近损害增长的临界点存在分歧。
  \item 欧元区内部财政分化（法意利差走阔）是否构成欧元系统性风险，不同机构看法不一。
\end{itemize}
\subsubsection*{机构观点}
\paragraph{MS} MS认为欧洲央行加息本身不会削弱欧元，除非市场开始定价增长受损或公共债务不确定性。当前欧元区1年期与5年5年远期利率之差约-78bp，远低于历史案例中货币明显转弱的阈值（100bp以上）。历史数据显示，G10货币在进入紧缩阶段后12个月平均相对回报为-0.8\%，但欧元在2023年7月进入紧缩后12个月仅跌0.7\%，表现优于均值。真正的风险在于欧元区内部利差走阔，若法国和意大利国债利差大幅扩大，将触发市场对财政风险的定价。
{\small\color{refgray}数据：欧元区1年期与5年5年远期利率之差约-78bp\par}
{\small\color{refgray}数据：紧缩程度达100bp以上的13个案例中，货币紧缩前12个月平均回报3.7\%，紧缩后12个月平均-0.8\%\par}
{\small\color{refgray}数据：欧元2023年7月进入紧缩后12个月仅跌0.7\%\par}
{\small\color{refgray}数据：新西兰1997年和2005年两次紧缩导致纽元大跌，因增长不确定性已浮现\par}
{\small\color{refgray}边际变化：此前市场普遍认为加息利好欧元，但MS指出需区分紧缩背景；当前欧元区利率水平尚未触发历史转弱阈值，但内部财政分化成为新关注点。\par}
\subsubsection*{关键数据}
\begin{itemize}
  \item 欧元区1年期与5年5年远期利率之差约-78bp
  \item 紧缩程度达100bp以上的13个案例中，货币紧缩前12个月平均回报3.7\%，紧缩后12个月平均-0.8\%
  \item 欧元2023年7月进入紧缩后12个月仅跌0.7\%
  \item 新西兰1997年和2005年两次紧缩导致纽元大跌
\end{itemize}
\begin{figure}[htbp]
\centering
\includegraphics[width=0.88\linewidth]{figures/fig_070.jpg}
\caption{Exhibit 1 - Exhibit 3: Long-end rate correlations are also positive, but more variable over time}
\end{figure}
\begin{figure}[htbp]
\centering
\includegraphics[width=0.88\linewidth]{figures/fig_073.jpg}
\caption{Exhibit 3 - Exhibit 3: Long-end rate correlations are also positive, but more variable over time Exhibit 4: JPY again screens as most correlated to long-end yield differentials These correlations have entered market focus in recent weeks a}
\end{figure}
\begin{figure}[htbp]
\centering
\includegraphics[width=0.88\linewidth]{figures/fig_071.jpg}
\caption{Exhibit 1 - Exhibit 3: Long-end rate correlations are also positive, but more variable over time Exhibit 4: JPY again screens as most correlated to long-end yield differentials}
\end{figure}
\begin{figure}[htbp]
\centering
\includegraphics[width=0.88\linewidth]{figures/fig_072.jpg}
\caption{Exhibit 2 - Exhibit 3: Long-end rate correlations are also positive, but more variable over time Exhibit 4: JPY again screens as most correlated to long-end yield differentials These correlations have entered market focus in recent weeks a}
\end{figure}
{\small\color{refgray}本节综合 1 篇研究；涉及 MS。}
\newpage
\subsection{科技-半导体与设备（1）：AI驱动需求结构性扩张，设备商受益分化}
\textbf{AI基础设施投资持续超预期，带动WFE市场2027年有望突破2000亿美元，但设备商受益程度因DRAM敞口、英特尔资本开支结构变化及光子学需求确认而显著分化。}
\subsubsection*{主流共识}
\begin{itemize}
  \item AI基础设施资本开支周期未见顶，Alphabet、英特尔等均上调指引，2027年WFE市场预计超2000亿美元。
  \item 光子学（Photonics-SOI）需求从远期故事变为可追踪增长，Soitec、Disco等确认AI数据中心光互连需求加速。
  \item CPU服务器需求受Agentic AI拉动被低估，通用服务器增长将支撑BMC等供应链。
\end{itemize}
\subsubsection*{分歧与条件差异}
\begin{itemize}
  \item DRAM资本开支节奏：最大DRAM制造商订单集中度导致日本设备商业绩分化，Ebara超预期而部分公司面临下调风险。
  \item 英特尔资本开支结构：厂房建设基本完成，新增支出主要流向设备采购，日本设备商（Lasertec、东京电子、Disco）受益程度高于美国同行。
  \item 台积电“超预期下跌”历史类比：NOM认为本次类似2021年假信号而非2022年周期见顶，但市场情绪仍存分歧。
  \item MLCC供需：三星电机提前锁定2027年产能，但市场对AI服务器MLCC增长倍数（6.4倍）的可持续性存在不同看法。
\end{itemize}
\subsubsection*{机构观点}
\paragraph{JPM} WFE市场2026-2028年预计分别达1590亿、2050亿、2370亿美元，年增幅16-29\%，但存在上调空间，因台积电资本支出上升、DRAM预期走高及Terafab概念推进。铠侠股价调整后，市场关注超大规模云厂商支出持续性及低成本AI模型对存储需求的影响，但杰文斯悖论（效率提升刺激总消耗）已缓解担忧。Soitec光子学业务FY27收入将环比增长，容量预留协议提升需求可见性，新加坡厂提前量产，RF-SOI库存正常化。
{\small\color{refgray}数据：WFE 2026E: 1590亿美元, 2027E: 2050亿美元, 2028E: 2370亿美元\par}
{\small\color{refgray}数据：铠侠FY27 EPS市场预期15000日元 vs 机构预期14149日元\par}
{\small\color{refgray}数据：Soitec F1Q27营收1.13亿欧元，超预期6\%\par}
{\small\color{refgray}数据：Soitec F2Q27指引同比>30\%，市场共识仅4\%\par}
{\small\color{refgray}边际变化：Soitec光子学从远期故事变为可追踪增长曲线，容量预留协议和新加坡厂提前量产提升收入可见性。\par}
\paragraph{MS} 云端半导体需求持续扩张，CPU与GPU服务器市场总量2027年前未见顶。CPU服务器在中国区超预期增长，受AMD多核、英伟达Vera捆绑及海光性能追赶驱动，AI通用服务器2027年将占CPU需求45\%。GPU服务器向超级吊舱架构演进，国产厂商展示64-128卡互联系统。Alphabet上调FY26资本支出指引至1950-2050亿美元，TPU系统2027年集中确认收入。日本设备商业绩分化，Ebara精密机械订单指引可能从4050亿日元上调至4500亿日元（同比+50\%），而部分公司面临下调风险。英特尔2027年WFE贡献约250亿美元，非美设备商（东京电子、Lasertec）受益明显。
{\small\color{refgray}数据：Alphabet FY26资本支出指引上调至1950-2050亿美元\par}
{\small\color{refgray}数据：Alphabet Q2资本支出449亿美元，环比+26\%，同比+100\%\par}
{\small\color{refgray}数据：Ebara精密机械订单计划可能从4050亿上调至4500亿日元\par}
{\small\color{refgray}数据：英特尔2027年WFE贡献约250亿美元\par}
{\small\color{refgray}边际变化：英特尔资本支出结构从厂房建设转向设备采购，2026年设备支出同比+40\%，非美设备商受益逻辑强化。\par}
\paragraph{NOM} 台积电7月业绩超预期后股价下跌7.3\%，类似情形过去十年仅出现四次。2022年4月下跌被证明是周期见顶，但本次由AI结构性需求驱动，与2021年假信号更相似。自2023年2月AI“iPhone时刻”以来，台积电、SOX和台湾OTC电子指数已累计上涨105\%、92\%和59\%，短期回调属正常整理，AI周期仍在持续。
{\small\color{refgray}数据：台积电业绩超预期后T+1日股价下跌7.3\%（台股）/2\%（ADR）\par}
{\small\color{refgray}数据：过去十年40次业绩会中仅4次出现“超预期下跌”\par}
{\small\color{refgray}数据：2022年4月下跌后台积电1个月跌11.9\%，SOX跌9.4\%\par}
{\small\color{refgray}数据：自2023年2月AI“iPhone时刻”至2026年7月，台积电累计上涨105\%\par}
{\small\color{refgray}边际变化：本次下跌与2022年周期见顶有本质区别，AI结构性需求提供支撑，短期整理而非趋势反转。\par}
\paragraph{GS} Disco Q1营业利润490亿日元符合预期，但Q2出货预测1410亿日元显著高于市场预期的1300亿日元，生成式AI和硅光子（CPO）设备出货扩大是主要推动力。英特尔Q2营收161亿美元超预期，DCAI收入63亿美元（同比+59\%），资本支出上调至200亿美元以上，设备支出同比+40\%，日本设备商（Lasertec、东京电子、Disco）受益明确。三星电机一个月内获两笔MLCC大单合计7490亿韩元，提前锁定2027年产能，AI服务器MLCC收入预计2025-2028年增长6.4倍。
{\small\color{refgray}数据：Disco Q2出货预测1410亿日元 vs 市场预期1300亿日元\par}
{\small\color{refgray}数据：英特尔Q2营收161亿美元，超市场预期11.7\%\par}
{\small\color{refgray}数据：英特尔DCAI收入63亿美元，同比+59\%，环比+24\%\par}
{\small\color{refgray}数据：三星电机两笔MLCC合同合计7490亿韩元，占去年MLCC收入16\%\par}
{\small\color{refgray}边际变化：英特尔资本支出结构转向设备采购，Disco CPO设备出货确认，三星电机MLCC大单提前锁定2027年产能，供需趋紧信号明确。\par}
\subsubsection*{关键数据}
\begin{itemize}
  \item WFE 2027E: 2020亿美元（MS），英特尔贡献约250亿美元
  \item Alphabet FY26资本支出指引上调至1950-2050亿美元
  \item Soitec F2Q27指引同比>30\%，市场共识仅4\%
  \item 英特尔2026年设备支出同比+40\%，全年资本支出超200亿美元
  \item 三星电机一个月内两笔MLCC大单合计7490亿韩元
  \item Disco Q2出货预测1410亿日元，高于市场预期
  \item 台积电业绩超预期后股价下跌7.3\%，历史第四次
  \item Ebara精密机械订单指引可能上调至4500亿日元（同比+50\%）
\end{itemize}
\begin{figure}[htbp]
\centering
\includegraphics[width=0.88\linewidth]{figures/fig_064.jpg}
\caption{Exhibit 1 - Exhibit 2: Aspeed expects its BMC TAM to reach 65.7mn units in 2030}
\end{figure}
\begin{figure}[htbp]
\centering
\includegraphics[width=0.88\linewidth]{figures/fig_137.jpg}
\caption{Exhibit 1 - Exhibit 1: Intel sales exposure across our Japan SPE coverage (FY25E GSe) Exhibit 2: Target prices, methodologies and risks for stocks mentioned \textbackslash{}*on APAC Conviction List}
\end{figure}
\begin{figure}[htbp]
\centering
\includegraphics[width=0.88\linewidth]{figures/fig_065.jpg}
\caption{Exhibit 2 - Exhibit 2: Aspeed expects its BMC TAM to reach 65.7mn units in 2030 \#\# BMC TAM Forecast Revised \textbackslash{}- We see huge demand for AI-related General Server for agentic AI workload and simple inference tasks. - BMC TAM is estimated to be 65}
\end{figure}
\begin{figure}[htbp]
\centering
\includegraphics[width=0.88\linewidth]{figures/fig_066.jpg}
\caption{Exhibit 3 - Exhibit 3: Aspeed's product roadmap Exhibit 4: Aspeed Product Roadmap Exhibit 5: Server-related companies maintain high Y/Y revenue growth – but slowing a bit because of a high base}
\end{figure}
{\small\color{refgray}本节综合 12 篇研究；涉及 JPM、MS、NOM、GS。}
\newpage
\subsection{功率半导体：短期利润承压，产品线扩张与长期毛利率修复成焦点}
\textbf{斯达半导2026年第二季度利润同比大幅下滑，主要受原材料成本上升和研发投入增加拖累，但环比改善且符合预期。市场关注点在于公司从IGBT向SiC、GaN等新产品的扩张能否在2027年后驱动毛利率回升，以及价格传导机制的落地节奏。}
\subsubsection*{主流共识}
\begin{itemize}
  \item 原材料成本上升是短期毛利率下行的主要压力，客户虽认可涨价但传导需时间。
  \item 公司正从IGBT向SiC MOSFET、GaN、DrMOS及MCU全面拓展，产品线扩张是长期增长核心。
  \item 2026年第二季度利润环比增长48\%，显示经营边际改善，但同比-77\%反映高基数效应。
\end{itemize}
\subsubsection*{分歧与条件差异}
\begin{itemize}
  \item 毛利率触底时间存在分歧：部分观点认为2026年下半年毛利率将企稳，另一部分认为需到2027年新品放量后。
  \item 研发投入强度：有分析认为当前高研发投入是必要战略，另一些担忧其持续压制利润。
\end{itemize}
\subsubsection*{机构观点}
\paragraph{GS} 斯达半导2026年第二季度利润中值3900万元，同比-77\%，环比+48\%，符合该机构预期但低于彭博一致预期57\%。利润下滑主因原材料成本上升致毛利率从2025年Q1的30.4\%降至2026年Q1的21.1\%，Q2预计进一步降至21.4\%；同时研发支出增加推高运营费用率。公司产品线从IGBT向SiC、GaN、DrMOS及MCU扩张，终端客户已认可成本上涨但调价需走流程，短期毛利率变化未完全明确。该机构下调2026-2028年毛利率预期1.0-1.5个百分点，上调运营费用率，预计长期毛利率随新品爬坡逐步改善至2028年的28.0\%。
{\small\color{refgray}数据：2026年Q2利润中值3900万元，同比-77\%，环比+48\%\par}
{\small\color{refgray}数据：毛利率从2025年Q1的30.4\%降至2026年Q1的21.1\%，Q2预计21.4\%\par}
{\small\color{refgray}数据：运营费用率从2025年的17.2\%升至2026年Q1的19.0\%\par}
{\small\color{refgray}数据：2026-2028年预测毛利率分别为25.2\%/27.8\%/28.0\%\par}
{\small\color{refgray}边际变化：相比此前预期，该机构下调了2026-2028年毛利率预测1.0-1.5个百分点，并上调运营费用率，反映短期成本压力超预期。\par}
\subsubsection*{关键数据}
\begin{itemize}
  \item 2026年Q2利润中值3900万元，同比-77\%，环比+48\%
  \item 毛利率从2025年Q1的30.4\%降至2026年Q1的21.1\%，Q2预计21.4\%
  \item 运营费用率从2025年的17.2\%升至2026年Q1的19.0\%
  \item 2026-2028年预测毛利率分别为25.2\%/27.8\%/28.0\%
  \item 2026年Q2利润较彭博一致预期低57\%
\end{itemize}
\begin{figure}[htbp]
\centering
\includegraphics[width=0.88\linewidth]{figures/fig_152.jpg}
\caption{Exhibit 3 - Exhibit 3: StarPower local and global peers' 2027 P/E and 27-28E earnings growth Exhibit 4: StarPower: 12m forward P/E}
\end{figure}
\begin{figure}[htbp]
\centering
\includegraphics[width=0.88\linewidth]{figures/fig_153.jpg}
\caption{Exhibit 3 - Exhibit 3: StarPower local and global peers' 2027 P/E and 27-28E earnings growth Exhibit 4: StarPower: 12m forward P/E}
\end{figure}
{\small\color{refgray}本节综合 1 篇研究；涉及 GS。}
\newpage
\subsection{AI与软件：云加速与模型分化}
\textbf{AI基础设施投资持续加码，云业务成为主要受益者，但企业软件支出出现结构性分化：AI原生需求驱动ServiceNow和微软Azure超预期，而IBM等传统IT厂商则面临客户预算临时性再分配的挑战。中国AI模型在编码能力上取得突破，加剧了全球AI竞争格局的不确定性。}
\subsubsection*{主流共识}
\begin{itemize}
  \item AI基础设施投资持续高增长，云业务（特别是Google Cloud和Azure）成为AI需求释放的主要载体，收入增速显著高于传统IT支出。
  \item 企业AI应用从概念验证进入规模化部署阶段，AI相关年度合同价值（ACV）和用户数快速增长，如ServiceNow AI ACV突破10亿美元，Gemini月活达9.5亿。
  \item AI芯片需求结构性分化，训练端CoWoS产能紧张边际缓解，但推理端（尤其是中国AI推理）成为新的增量来源。
\end{itemize}
\subsubsection*{分歧与条件差异}
\begin{itemize}
  \item 企业IT支出优先级：IBM客户将预算从大型机转向分布式基础设施，是临时性资金再分配还是系统性优先级重置？MS认为管理层预期下半年回升逻辑成立但容错空间有限。
  \item AI模型竞争格局：Kimi K3等中国开源模型在编码能力上接近美国前沿水平，是否会加速闭源模型商业化进程放缓？NOM认为长期推动行业增长，但短期可能加剧竞争。
  \item AI芯片需求结构：云端训练端CoWoS产能紧张缓解，但推理端需求能否持续填补产能缺口？MS指出高端封装产能释放慢于芯片设计复杂度提升。
\end{itemize}
\subsubsection*{机构观点}
\paragraph{Bernstein} Alphabet 2Q26搜索收入符合预期（+17\%），云业务大幅超预期（+82\%至248亿美元），但市场关注模型进展和现金流转负（-59亿美元）。AI概览和AI模式月活突破10亿，但搜索增速下半年将放缓至中双位数。云业务利润率36\%，但租用第三方算力可能影响利润率。
{\small\color{refgray}数据：Alphabet 2Q26收入同比+24\%至1198亿美元\par}
{\small\color{refgray}数据：云业务收入+82\%至248亿美元，超预期10\%\par}
{\small\color{refgray}数据：运营利润率约34\%\par}
{\small\color{refgray}数据：现金流为负59亿美元\par}
{\small\color{refgray}边际变化：首次出现季度现金流转负，云业务超预期但利润率面临调整压力\par}
\paragraph{Bernstein} ServiceNow 2Q26订阅收入超指引230bps，一半来自净新增ACV，一半来自本地部署提前续约（AI控制塔和安全产品驱动）。Q3 cRPO指引环比增长强劲，但提前续约拉低Q3收入约100bps。管理层对下半年保守但需求良好。
{\small\color{refgray}数据：订阅收入超指引230bps\par}
{\small\color{refgray}数据：AI控制塔和安全产品驱动提前续约\par}
{\small\color{refgray}数据：Q3 cRPO指引环比增长强劲\par}
{\small\color{refgray}数据：提前续约拉低Q3收入约100bps\par}
{\small\color{refgray}边际变化：提前续约是时间平移而非增量，净新增ACV才是真正的超预期来源\par}
\paragraph{Bernstein} IBM FQ2 2026收入172亿美元（同比+1\%），低于共识。大型机收入下降42\%，但分布式基础设施增长37\%（创纪录）。客户因内存短缺和涨价预期将资金转向服务器和存储，属于临时性资金再分配，非需求结构变化。
{\small\color{refgray}数据：IBM FQ2收入172亿美元，同比+1\%\par}
{\small\color{refgray}数据：大型机收入下降42\%\par}
{\small\color{refgray}数据：分布式基础设施增长37\%\par}
{\small\color{refgray}数据：全球70\%交易价值仍运行在IBM Z上\par}
{\small\color{refgray}边际变化：客户预算因内存涨价预期临时调整，大型机订单后移但未流失\par}
\paragraph{JPM} ServiceNow F2Q26 AI ACV突破10亿美元，cRPO增速21.5\%（固定汇率），高于预期约200bps。AI控制塔成为企业AI工作流核心入口，首次购买Agentic AI客户数同比+45\%，Pro Plus SKU提价超30\%。AI净新增ACV环比+40\%，含5个以上AI产品的交易量同比+5.5倍。
{\small\color{refgray}数据：AI ACV突破10亿美元\par}
{\small\color{refgray}数据：cRPO增速21.5\%，高于预期200bps\par}
{\small\color{refgray}数据：首次购买Agentic AI客户数同比+45\%\par}
{\small\color{refgray}数据：Pro Plus SKU提价超30\%\par}
{\small\color{refgray}边际变化：AI从概念增量变为实质拉动客户升级和扩购，缓解了市场对AI替代传统软件的担忧\par}
\paragraph{MS} IBM 2Q26营收172亿美元低于预期，但成本优化使税前利润率同比扩张30bps。管理层预期下半年回升，但需3Q/4Q软件环比增速超过去5年任何时期。约1/3延迟的ELA已在3Q前三周完成签约，快于正常节奏。
{\small\color{refgray}数据：IBM 2Q26营收172亿美元，低于预期约10亿美元\par}
{\small\color{refgray}数据：税前利润率同比扩张30bps\par}
{\small\color{refgray}数据：全年PTI利润率指引上调至同比+100bps\par}
{\small\color{refgray}数据：约1/3延迟的ELA在3Q前三周完成签约\par}
{\small\color{refgray}边际变化：成本效率是当前最可靠的回报支撑，但下半年增长目标要求历史级别的环比增速\par}
\paragraph{MS} 微软FY4Q26 Azure固定汇率增速有望达41\%，高于指引上限1个百分点，受渠道调研、GPU产能改善和CIO调查三重支撑。Copilot付费席位净增可能超600万，E7 SKU和按用量计费叠加成为最大ARPU提升机会。
{\small\color{refgray}数据：Azure FY4Q26增速预计41\%（固定汇率）\par}
{\small\color{refgray}数据：Copilot付费用户超2000万，净增可能超600万\par}
{\small\color{refgray}数据：CIO调查中微软支出预期增速+7.6\%\par}
{\small\color{refgray}数据：Azure AI增速约100\%同比\par}
{\small\color{refgray}边际变化：Azure和Copilot进入同步加速阶段，市场尚未充分定价\par}
\paragraph{MS} AI芯片需求结构性分化：CoWoS产能紧张边际缓解，但高端封装产能释放慢于芯片设计复杂度提升。中国AI推理需求成为新变量，DeepSeek激发本土算力采购意愿。非AI芯片产能被挤压和成本上升改变供应链优先级。
{\small\color{refgray}数据：CoWoS产能紧张边际缓解\par}
{\small\color{refgray}数据：高端封装产能释放慢于芯片设计复杂度提升\par}
{\small\color{refgray}数据：中国AI推理需求成为新增量\par}
{\small\color{refgray}数据：非AI芯片产能被挤压\par}
{\small\color{refgray}边际变化：CoWoS瓶颈从产能总量转向特定层数和线宽组合，中国推理需求成为增量变量\par}
\paragraph{NOM} xAI对GB300订单大幅上修至超13k机架（此前预估6-8k），缓解了GB300向VR200过渡的出货衔接担忧。2026年GB/VR机架出货预测从54.5k上调至62k台。Supermicro 4QFY26毛利率达15-17\%，远高于此前预期，受益于有利客户与产品组合。
{\small\color{refgray}数据：xAI GB300订单超13k机架\par}
{\small\color{refgray}数据：2026年GB/VR机架出货预测上调至62k台\par}
{\small\color{refgray}数据：Supermicro 4QFY26毛利率15-17\%\par}
{\small\color{refgray}数据：SpaceX 2026年6月公开市场融资857亿美元\par}
{\small\color{refgray}边际变化：Neocloud购买力超预期，用实际订单填补了产品代际切换的衔接段\par}
\paragraph{NOM} Kimi K3（2.8万亿参数混合专家模型）将Transformer与循环模型结合，编码能力接近美国前沿水平。KDA循环模型降低KV缓存需求，MLA层保留细粒度检索能力。长期推动AI行业增长，但短期可能延缓美国闭源模型商业化进程。
{\small\color{refgray}数据：Kimi K3参数规模2.8万亿\par}
{\small\color{refgray}数据：896个专家模块，每token激活16个\par}
{\small\color{refgray}数据：最大上下文长度100万token\par}
{\small\color{refgray}数据：编码能力接近美国前沿水平\par}
{\small\color{refgray}边际变化：循环模型落地时间点比预期更早，为中国AI模型开发提供新路径\par}
\paragraph{NOM} Alphabet 2Q26资本支出449亿美元（同比+100\%），全年计划上调至1950-2050亿美元，2027年将继续显著增长。云业务收入248亿美元（+82\%），积压订单5140亿美元。Gemini月活9.5亿，模型API每分钟处理220亿token。自研TPU首次交付外部客户。
{\small\color{refgray}数据：Alphabet 2Q26资本支出449亿美元，同比+100\%\par}
{\small\color{refgray}数据：全年资本支出计划上调至1950-2050亿美元\par}
{\small\color{refgray}数据：云业务积压订单5140亿美元\par}
{\small\color{refgray}数据：Gemini月活9.5亿\par}
{\small\color{refgray}边际变化：资本支出再次提速，2027年将继续显著增长，供应链受益明确\par}
\paragraph{GS} AMD与Anthropic达成战略合作，Anthropic将从2027年上半年部署2GW Instinct MI450 GPU，AMD计划进行最高50亿美元战略股权研究。微软将从2026年下半年部署Helios机架用于Azure推理。AMD将2030年数据中心AI加速器市场预期从2000亿上调至1.4万亿美元。
{\small\color{refgray}数据：Anthropic部署2GW MI450 GPU\par}
{\small\color{refgray}数据：AMD计划对Anthropic进行50亿美元战略股权研究\par}
{\small\color{refgray}数据：微软2026年下半年部署Helios机架\par}
{\small\color{refgray}数据：2030年数据中心AI加速器市场预期上调至1.4万亿美元\par}
{\small\color{refgray}边际变化：AMD同时获得前沿AI模型公司和云服务商客户，50亿美元股权研究加深绑定\par}
\paragraph{Bernstein} GE Aerospace 在 7 月 16 日发布的二季度财报中，不仅再次超出市场预估，还上调了全年指引。与一季度单纯超预期不同，这次管理层对全年展望的信心明显增强。报告认为，尽管伊朗局势对燃油价格的影响仍存在不确定性，但最糟糕的阶段可能已经过去，宏观宏观环境保持稳健——这对航空发动机售后市场而言是最关键的支撑因素。
{\small\color{refgray}数据：二季度 GE Aerospace 非 GAAP 每股表现为 2.02 美元，高于市场预估的 1.86 美元。调整后收入 126 亿美元，同样超出共识预估的 119 亿美元。商用发动机与服务（CES）部门收入 97.3 亿美元，高于预估的 91.6 亿美元；国防与推进技术（DP\&T）部门收入 34.4 亿美元，高于预估的 32.0 亿美元。两个部门的利润均超出...\par}
{\small\color{refgray}数据：全年指引的核心变化来自商用发动机服务。GE 将 2026 年 CES 部门收入增长预期从此前的“中十几个百分点”上调至“约 20\%”，其中服务收入增长预期从“中十几个百分点”上调至“低二十几个百分点”。Bernstein估算，这意味着商用服务收入将增加约 10 亿美元以上。\par}
\subsubsection*{关键数据}
\begin{itemize}
  \item Alphabet 2Q26资本支出449亿美元，同比+100\%
  \item ServiceNow AI ACV突破10亿美元，cRPO增速21.5\%
  \item xAI GB300订单超13k机架，2026年GB/VR机架出货预测上调至62k台
  \item Kimi K3参数规模2.8万亿，编码能力接近美国前沿水平
  \item Azure FY4Q26增速预计41\%（固定汇率）
  \item Google Cloud积压订单5140亿美元
  \item IBM大型机收入下降42\%，分布式基础设施增长37\%
  \item AMD将2030年数据中心AI加速器市场预期上调至1.4万亿美元
\end{itemize}
\begin{figure}[htbp]
\centering
\includegraphics[width=0.88\linewidth]{figures/fig_001.jpg}
\caption{EXHIBIT 4 - EXHIBIT 6: Digital ads growth (\%) Reported ad revenue Y/Y growth (\%) through 2Q26}
\end{figure}
\begin{figure}[htbp]
\centering
\includegraphics[width=0.88\linewidth]{figures/fig_016.jpg}
\caption{EXHIBIT 1 - EXHIBIT 1: The guided QoQ increase in cRPO for Q3 is strong compared to previous years. EXHIBIT 2: The YoY growth in incremental Q3 cRPO laps a tough -1 year comp, but as is seen in ex 1 the \textbackslash{}\$-level is very strong.}
\end{figure}
\begin{figure}[htbp]
\centering
\includegraphics[width=0.88\linewidth]{figures/fig_075.jpg}
\caption{Exhibit 9 - Exhibit 3: Microsoft and Amazon Continue to Lead as the Top Incremental Share Gainers due to Cloud Shift in 2026...}
\end{figure}
\begin{figure}[htbp]
\centering
\includegraphics[width=0.88\linewidth]{figures/fig_069.jpg}
\caption{Exhibit 1 - Exhibit 1: Our base case of 11.6x P/FCF reflects a regression of software and services peers, and leads to our PT of \textbackslash{}\$190 P/FCFe 2026 vs. FCF Growth What we learned from IBM's 2Q26 earnings conference call and IR callback: That}
\end{figure}
{\small\color{refgray}本节综合 12 篇研究；涉及 Bernstein、JPM、MS、NOM、GS。}
\newpage
\subsection{AI商业化进入分化期：云与软件验证收入，但利润稀释与硬件分流风险并存}
\textbf{AI投资正从基础设施层向平台和应用层迁移，Google Cloud和SAP的云积压订单验证了AI商业化的收入兑现能力，但IBM的硬件分流和SAP的利润稀释表明，AI价值分配的不确定性正在加剧，市场需区分总量增长与分配风险。}
\subsubsection*{主流共识}
\begin{itemize}
  \item AI商业化正从基础设施层向平台和应用层迁移，云业务收入加速是核心验证信号。
  \item 资本支出持续扩张，但市场更关注其能否转化为可量化的客户合同和收入增长。
  \item AI模型竞争从单一性能转向“性能+成本”双维度，性价比成为差异化关键。
\end{itemize}
\subsubsection*{分歧与条件差异}
\begin{itemize}
  \item AI价值分配：GS认为分配性事件（如谷歌融资、Meta建云）对大盘影响有限，但IBM的硬件分流表明，客户预算从软件向硬件转移可能改变特定公司的收入质量。
  \item 利润与增长平衡：SAP云合同连续超预期但利润指引下调，市场对AI投入的利润稀释效应存在分歧——是短期阵痛还是长期结构性问题。
  \item 中国AI模型路径：MiniMax强调自建算力的成本优势，而其他中国模型厂商可能更依赖外部算力租赁，成本结构差异将影响竞争格局。
  \item AI对软件行业的影响：GS认为AI嵌入交易（如SAP 90\%大额交易含AI）提升收入可见度，但IBM的交易处理业务因硬件分流而下滑，AI对传统软件是赋能还是替代存在分歧。
\end{itemize}
\subsubsection*{机构观点}
\paragraph{GS} Alphabet在AI领域的分发与算力规模优势正转化为可量化的收入：Google Cloud 2Q26同比增长82\%，积压订单达5140亿美元，近90\%财富100强使用Gemini Enterprise。资本支出指引上调至1950-2050亿美元（2026年），2027年预计3500亿美元，资本密集度将从2025年的26.7\%升至2027年的63.8\%。GS认为市场低估了其AI商业化的长期价值，并上调2026/2027年云计算收入预期至4330亿/5488亿美元。
{\small\color{refgray}数据：Google Cloud 2Q26同比增长82\%\par}
{\small\color{refgray}数据：积压订单环比增加约500亿美元至5140亿美元\par}
{\small\color{refgray}数据：近90\%财富100强使用Gemini Enterprise\par}
{\small\color{refgray}数据：2026年资本支出指引上调至1950-2050亿美元\par}
{\small\color{refgray}边际变化：此前市场关注AI资本支出的烧钱风险，但GS强调资本支出扩张与云计算收入加速的直接关联，且直接TPU硬件销售/承诺占比上升，算力投资正转变为可计量的客户合同。\par}
\paragraph{GS} AI风险需区分总量影响（改变AI创新总价值预期）与分配影响（改变价值在不同主体间的分配）。DeepSeek事件（2025年1月）和2026年2月的广泛AI担忧属于总量风险，导致标普500下跌1.5\%-1.7\%；而谷歌股权融资（2026年6月）和Meta建云（2026年7月）等分配性事件对大盘影响极小（<0.2\%）。分配风险对特定板块冲击更大，但管理难度更高。
{\small\color{refgray}数据：DeepSeek事件：标普500下跌1.5\%\par}
{\small\color{refgray}数据：2026年2月AI担忧：标普500下跌1.7\%\par}
{\small\color{refgray}数据：谷歌股权融资：标普500变动<0.2\%\par}
{\small\color{refgray}数据：Meta建云：标普500变动<0.2\%\par}
{\small\color{refgray}边际变化：此前市场将AI相关波动视为整体风险，GS首次系统区分总量与分配风险，指出分配性事件对大盘影响有限但板块冲击更大，为组合管理提供新框架。\par}
\paragraph{GS} IBM面临客户需求从软件向硬件的结构性分流：2Q26收入172亿美元（同比+1\%），软件收入77.6亿美元（同比+5\%，低于预期11\%），交易处理业务下降9\%；基础设施收入38.4亿美元（同比-7.4\%），大型机减少42\%。2026年软件增速预期从>10\%下调至6\%-8\%，收入组合更依赖硬件，公司需更长时间重建市场信任。
{\small\color{refgray}数据：2Q26总收入172亿美元，同比+1\%\par}
{\small\color{refgray}数据：软件收入77.6亿美元，同比+5\%，低于预期11\%\par}
{\small\color{refgray}数据：交易处理业务下降9\%\par}
{\small\color{refgray}数据：基础设施收入38.4亿美元，同比-7.4\%\par}
{\small\color{refgray}边际变化：此前市场预期IBM软件增长稳定，但客户在内存和组件价格上涨背景下优先采购硬件，导致软件收入低于预期且交易处理业务持续调整，收入质量下降。\par}
\paragraph{GS} SAP云业务前瞻指标强劲但利润端承压：2Q26当前云合同（CCB）达229亿欧元（固定汇率同比+26\%），连续两个季度超预期，Cloud ERP Suite驱动，超过90\%的50大交易嵌入AI和SAP Business Data Cloud。但非IFRS营业利润同比仅+5\%（低于预期8\%），全年营业利润指引因收购稀释下调1亿欧元至118-122亿欧元。
{\small\color{refgray}数据：CCB达229.29亿欧元，固定汇率同比+26\%\par}
{\small\color{refgray}数据：CCB有机增长约24.7\%，高于GS预期24.8\%和一致预期24.3\%\par}
{\small\color{refgray}数据：云收入62.8亿欧元，同比+24\%\par}
{\small\color{refgray}数据：非IFRS营业利润32.1亿欧元，同比+5\%，低于预期8\%\par}
{\small\color{refgray}边际变化：此前市场关注云收入增长，但SAP利润指引下调反映收购、股权激励和研发投入的稀释效应，市场需重新评估收入韧性与利润稀释的平衡。\par}
\paragraph{GS} MiniMax通过自建算力的时间窗口形成成本优势：自2025年9月起大规模自营算力（购置+长租），当时GPU租金低于当前水平。M3模型日处理token达万亿级别，API已实现毛利率。即将发布的M3 Pro（约3万亿参数，2026年9-10月）通过稀疏注意力、激活参数优化和KV缓存技术，推理成本有望明显低于同类模型。中国大模型竞争从单一性能转向“价格+性能”双维度。
{\small\color{refgray}数据：M3模型日处理token达万亿级别\par}
{\small\color{refgray}数据：M3 Pro约3万亿参数级别，预计2026年9-10月发布\par}
{\small\color{refgray}数据：自2025年9月起大规模增加自运营算力规模\par}
{\small\color{refgray}数据：当时GPU长期租赁价格远低于当前水平\par}
{\small\color{refgray}边际变化：此前市场关注中国大模型的性能竞赛，但MiniMax强调性价比策略，自建算力的低成本优势（低点锁定长期资源）成为差异化竞争的关键。\par}
\subsubsection*{关键数据}
\begin{itemize}
  \item Google Cloud 2Q26同比增长82\%，积压订单5140亿美元
  \item Alphabet 2026年资本支出指引上调至1950-2050亿美元，2027年预计3500亿美元
  \item DeepSeek事件导致标普500下跌1.5\%，GS美国AI广泛篮子下跌9.5\%
  \item IBM软件收入同比+5\%，低于预期11\%，交易处理业务下降9\%
  \item SAP CCB达229亿欧元，固定汇率同比+26\%，连续两个季度超预期
  \item SAP非IFRS营业利润同比+5\%，低于预期8\%，全年指引下调1亿欧元
  \item MiniMax M3模型日处理token达万亿级别，M3 Pro约3万亿参数
  \item 超过90\%的SAP 50大交易嵌入AI和SAP Business Data Cloud
\end{itemize}
\begin{figure}[htbp]
\centering
\includegraphics[width=0.88\linewidth]{figures/fig_130.jpg}
\caption{Exhibit 1 - Exhibit 2: Our dominant focus has been on comparing value built into US companies and the potential capital revenues from AI Change in Market Cap/Valuation From Nov. 30 2022 to Jul. 22 2026 vs. GS Estimates of PDV of Potential AI C}
\end{figure}
\begin{figure}[htbp]
\centering
\includegraphics[width=0.88\linewidth]{figures/fig_131.jpg}
\caption{Exhibit 2 - Exhibit 2: Our dominant focus has been on comparing value built into US companies and the potential capital revenues from AI Change in Market Cap/Valuation From Nov. 30 2022 to Jul. 22 2026 vs. GS Estimates of PDV of Potential AI C}
\end{figure}
\begin{figure}[htbp]
\centering
\includegraphics[width=0.88\linewidth]{figures/fig_132.jpg}
\caption{Exhibit 3 - Exhibit 3: Views on the value of AI will shift if investors adjust their views on the key parameters Estimates of the PDV of AI revenues to US companies are sensitive to assumptions about each of these parameters, as we have illu}
\end{figure}
\begin{figure}[htbp]
\centering
\includegraphics[width=0.88\linewidth]{figures/fig_133.jpg}
\caption{Exhibit 6 - Exhibit 6: Single stock volatility has risen sharply over the last year while implied correlation has hit new lows Exhibit 7: High single stock volatility, low implied correlation reflects a focus on distributional issues \#\# Ag}
\end{figure}
{\small\color{refgray}本节综合 5 篇研究；涉及 GS。}
\newpage
\subsection{AI能源基础设施：供需再平衡与区域分化}
\textbf{全球数据中心与电力基础设施正经历结构性扩张，但供给端加速与需求可见度之间的差距正在制造新的分化：上游电力设备订单持续强劲，而下游数据中心运营商面临宏观估值压制与项目执行风险。}
\subsubsection*{主流共识}
\begin{itemize}
  \item 全球数据中心容量到2030年将翻倍以上，美国占净增量的60-70\%，对应约5.8万亿美元资本开支，超大规模云厂商年均约1万亿美元的支出可覆盖。
  \item 燃气轮机与电气化设备订单周期尚未见顶，GE Vernova等OEM积压订单延伸至2027-2028年，定价环境保持稳定。
  \item AI数据中心成为韩国电信运营商长期分化的核心驱动力，KT、SKT、LGU+均公布大规模AIDC扩容计划。
  \item 数据中心空置率持续压缩，服务商定价能力增强，美国主要市场每千瓦费用从2021年的约70美元升至接近250美元。
\end{itemize}
\subsubsection*{分歧与条件差异}
\begin{itemize}
  \item 燃气轮机订单是否接近峰值：JPM认为订单结构仍有支撑（正式订单为季度出货量3.7倍），而市场担忧2030年供给过剩。
  \item SKT的AIDC战略：GS认为当前股价已计入过多乐观预期，项目执行和回报可见度低；而JPM将SK Innovation的转型视为积极信号。
  \item 数据中心运营商估值：MS认为VNET基本面稳健但受离岸利率环境压制，而GS上调全球容量预期但强调供需仍紧张。
  \item 韩国电信公司AIDC路径：KT采用轻资产预销售模式，SKT采用先投入后签约模式，两者风险收益特征不同。
\end{itemize}
\subsubsection*{机构观点}
\paragraph{JPM} GE Vernova 2Q26业绩显示电气化板块收入同比增64\%（有机29\%），EBITDA利润率扩张至18.4\%，管理层上调全年收入指引至145-150亿美元。燃气轮机订单结构健康：新增20GW中18GW为预留协议，但10GW已转为正式订单，实际正式订单12GW为季度出货量3.7倍，积压订单从44GW增至53GW。市场对订单见顶担忧被放大，供需基本面在本十年末前仍有支撑。
{\small\color{refgray}数据：电气化收入同比+64\%（有机+29\%）\par}
{\small\color{refgray}数据：EBITDA利润率同比+4ppt至18.4\%\par}
{\small\color{refgray}数据：燃气轮机正式订单12GW，为季度出货量3.7倍\par}
{\small\color{refgray}数据：积压订单44→53GW，槽位预留56→63GW\par}
{\small\color{refgray}边际变化：市场此前担忧燃气轮机订单见顶，但JPM拆解后指出实际正式订单强劲，且电气化板块利润率持续提升，订单收入比1.7倍，表明高压输电设备需求周期仍在延续。\par}
\paragraph{JPM} SK Innovation出售SK China 20\%股权回笼9000亿韩元，战略重心从电动车电池转向AI能源基础设施。电池业务瘦身完成后，公司正成为SK集团AI全栈战略的能源供应商，已与SK Telecom合作在蔚山建设103MW AI数据中心（规划扩展至1GW），毗邻LNG联合循环发电厂。市场仍锚定旧叙事，但公司资产负债表趋于稳定，资本开支回落至2022年前水平。
{\small\color{refgray}数据：出售SK China 20\%股权回笼9000亿韩元\par}
{\small\color{refgray}数据：蔚山AI数据中心初始103MW，规划扩展至1GW\par}
{\small\color{refgray}数据：韩国Honam半导体集群电力需求超6GW\par}
{\small\color{refgray}数据：资本开支回落至2022年前水平\par}
{\small\color{refgray}边际变化：市场此前关注SK Innovation的电池亏损，但公司正转型为AI能源基础设施供应商，出售非核心资产回笼资金用于电力基础设施投资，角色转换尚未被充分定价。\par}
\paragraph{MS} VNET基本面稳健但市场表现持续承压，主要受离岸利率环境压制而非公司运营问题。预计2Q26营收同比+10.6\%，调整后EBITDA同比+23\%，但环比仅+1.1\%，主因芯片供应后置导致产能释放集中在下半年。山东高速与CATL交易已获批，超大规模客户新订单有望在业绩中公布。当前EV/EBITDA约10.5倍，远期2027年约9.6倍，DCF模型WACC假设8.7\%。
{\small\color{refgray}数据：2Q26营收预期26.91亿元，同比+10.6\%\par}
{\small\color{refgray}数据：调整后EBITDA预期9.02亿元，同比+23\%，环比+1.1\%\par}
{\small\color{refgray}数据：EBITDA利润率预期33.5\%\par}
{\small\color{refgray}数据：当前EV/EBITDA 10.5倍，远期2027年9.6倍\par}
{\small\color{refgray}边际变化：市场将VNET股价弱势归因于基本面，但MS指出核心压制因素为离岸利率环境，公司负债主要在境内融资，情绪面影响大于实质。芯片供应后置意味着下半年产能释放可能加速。\par}
\paragraph{GS} 全球数据中心容量预测上调：2030年供应预期从168GW升至217GW，较2025年净增116GW，对应约5.8万亿美元建设投入。美国市场2030年预计达125GW，占净增量60-70\%。尽管供给加速，空置率仍持续压缩，服务商定价能力增强（每千瓦费用从2021年约70美元升至接近250美元）。超大规模云厂商年均约1万亿美元资本开支可覆盖建设投入，且该共识未计入SpaceX AI、私营AI实验室等支出。
{\small\color{refgray}数据：2030年全球数据中心供应预期217GW（此前168GW）\par}
{\small\color{refgray}数据：2025年供应101GW（此前74GW）\par}
{\small\color{refgray}数据：净增116GW对应约5.8万亿美元建设投入\par}
{\small\color{refgray}数据：美国2030年预计125GW，占净增量60-70\%\par}
{\small\color{refgray}边际变化：GS大幅上调全球数据中心容量预测（较上一季度+5\%，较两季度前+13\%），但强调供需仍紧张，空置率持续压缩，服务商定价能力增强。新增容量中70-90\%来自零售/批发供应商和超大规模云厂商。\par}
\paragraph{GS} SKT设立AI数据中心子公司SK Hyper并承诺至2030年投入最高7500亿韩元股权资金，首期3300亿韩元7月底到位。但GS认为这并未改变战略不确定性轮廓：资金主要用于开发平台和前期基础设施（土地、电网、许可），而非5GW总规划的建设成本。项目层面关键变量——租约、融资、回报表现——仍缺乏可验证信息。当前股价99,500韩元已计入较多乐观预期。
{\small\color{refgray}数据：SK Hyper至2030年投入最高7500亿韩元\par}
{\small\color{refgray}数据：首期3300亿韩元7月底到位\par}
{\small\color{refgray}数据：首批约2GW项目位于蔚山和忠清西部\par}
{\small\color{refgray}数据：SKT此前公布5GW产能计划，2029年起分阶段投运\par}
{\small\color{refgray}边际变化：市场将SKT的AIDC子公司设立视为积极信号，但GS指出资金主要用于前期开发而非建设成本，且先投入后签约模式可能延长开发周期、增加额外资金需求，执行风险被低估。\par}
\subsubsection*{关键数据}
\begin{itemize}
  \item 2030年全球数据中心供应预期217GW（此前168GW），较2025年净增116GW
  \item 美国2030年数据中心容量预计125GW，占全球净增量60-70\%
  \item 每千瓦费用从2021年约70美元升至接近250美元
  \item GE Vernova电气化收入同比+64\%（有机+29\%），EBITDA利润率18.4\%
  \item 燃气轮机正式订单12GW，为季度出货量3.7倍
  \item SK Innovation出售SK China 20\%股权回笼9000亿韩元
  \item SKT承诺至2030年投入7500亿韩元于AI数据中心平台
  \item VNET 2Q26调整后EBITDA同比+23\%，环比仅+1.1\%
\end{itemize}
\begin{figure}[htbp]
\centering
\includegraphics[width=0.88\linewidth]{figures/fig_026.jpg}
\caption{Figure 1 - Figure 1: Backlog + SRAs GW ■ Firm Order Backlog ■ Slot Reservation Agreements}
\end{figure}
\begin{figure}[htbp]
\centering
\includegraphics[width=0.88\linewidth]{figures/fig_120.jpg}
\caption{Exhibit 1 - Exhibit 1: Global Data Center Capacity should scale to \textbackslash{}\textasciitilde{}217 MW by 2030 (+116 GW v. 2025) Gloabl Data Center Capacity (MWs) \& year-over-year change (\%) Global Data Center Capacity forecast estimates over time (MW) Exhibit 3: Gl}
\end{figure}
\begin{figure}[htbp]
\centering
\includegraphics[width=0.88\linewidth]{figures/fig_122.jpg}
\caption{Exhibit 2 - Exhibit 5: US data center capacity should grow to \textbackslash{}\textasciitilde{}125 GW by 2030 (+76 GW from 2025)}
\end{figure}
\begin{figure}[htbp]
\centering
\includegraphics[width=0.88\linewidth]{figures/fig_085.jpg}
\caption{Exhibit 1 - Exhibit 1: VNET - 2Q26 preview Exhibit 2: VNET – EV/EBITDA vs. US 10-year TIPS (inverse) – key valuation drag Asia Summer School 2026}
\end{figure}
{\small\color{refgray}本节综合 7 篇研究；涉及 JPM、MS、GS。}
\newpage
\subsection{MLCC与CCL涨价周期延续，面板定价承压分化}
\textbf{MLCC与覆铜板（CCL）涨价逻辑未变，产能偏紧与AI需求持续支撑中期叙事；但面板端因存储芯片涨价挤压客户预算，ASP指引弱于预期，产业链利润分化明显。}
\subsubsection*{主流共识}
\begin{itemize}
  \item MLCC与CCL产能偏紧，涨价周期至少延续至2027年，AI需求是核心驱动。
  \item 上游电子布/玻璃布利润增速显著高于CCL和PCB，产业链利润从上游向下游递减。
  \item AI服务器对高容MLCC和高速CCL的需求持续提升，推动产品结构升级。
  \item 消费电子需求平稳，但AI相关拉货和客户预防性备货支撑短期出货。
\end{itemize}
\subsubsection*{分歧与条件差异}
\begin{itemize}
  \item MLCC涨价预期是否已充分反映：MS认为中期逻辑未变，但部分客户已开始预防性备货，可能提前透支需求。
  \item CCL涨价持续性分歧：Citi认为电子布供应偏紧至2027年底，价格仍有上行空间；但需关注低端玻璃布供给2027年底缓解后对价格的压制。
  \item 面板定价分化：GS认为存储芯片涨价挤压客户预算，导致OLED ASP指引弱于预期；但大尺寸OLED受益于赛事和电竞需求，利润率改善。
\end{itemize}
\subsubsection*{机构观点}
\paragraph{MS} MLCC中期叙事未变，产能仍需匹配AI高容需求。二季度营收442亿新台币（环比+16\%，同比+35\%），超预期驱动来自定价环境改善及客户预防性备货。三季度预计营收环比+8\%，毛利率40.3\%（环比+150bp），营业利润率28.9\%（环比+210bp），定价环境是核心驱动。营业利润率存在20bp上行空间。
{\small\color{refgray}数据：2Q26营收442亿新台币，环比+16\%，同比+35\%\par}
{\small\color{refgray}数据：2Q26毛利率38.8\%，环比+70bp；营业利润率26.8\%，环比+170bp\par}
{\small\color{refgray}数据：3Q26预计营收环比+8\%，毛利率40.3\%，营业利润率28.9\%\par}
{\small\color{refgray}数据：营业利润率存在20bp上行空间\par}
{\small\color{refgray}边际变化：二季度收入超预期并非来自消费电子回暖，而是AI服务器拉货叠加客户预防性备货；中低端MLCC供应预期变化显示高容产能挤出效应向下传导。\par}
\paragraph{Citi} CCL涨价周期仍有空间，标准无卤板报价已超2021年高点（当前>300元/张，2021年高点约280元）。管理层将标准品最低毛利率门槛从1Q26的15\%上调至3Q26的23-25\%，连续上调反映定价能力增强。电子布供应偏紧至少持续至2027年底，AI级电子布占产能20\%，2027-2028年扩产主要投向AI-CCL。产业链利润从上游向下游递减，电子布厂利润增速最高。
{\small\color{refgray}数据：标准无卤板报价>300元/张，2025年底160-170元，2021年高点约280元\par}
{\small\color{refgray}数据：标准品最低毛利率门槛：1Q26 15\% → 2Q26 18-20\% → 3Q26 23-25\%\par}
{\small\color{refgray}数据：电子布供应偏紧至少持续至2027年底\par}
{\small\color{refgray}数据：1H26净利润预告4.2-5亿元，同比+382-473\%\par}
{\small\color{refgray}边际变化：管理层对电子布供应偏紧的判断从2027年中延长至2027年底，比Citi此前预期更积极；CCL价格从2025年底160-170元升至当前300元，管理层预计年底达360-400元。\par}
\paragraph{GS} 三菱电机与索尼成立AI视觉传感器合资公司（2026年10月运营，三菱持股60\%），填补FA数字化战略中传感器层缺失。合资公司开发边缘AI+图像传感器技术，实现产线无人工干预监控。该合作补上三菱电机在PLC、伺服、机器人之外的“眼睛”环节，是工业视觉领域趋势。
{\small\color{refgray}数据：合资公司Advanced Vision Solutions，三菱电机持股60\%，索尼半导体持股40\%\par}
{\small\color{refgray}数据：2026年10月开始运营，纳入三菱电机合并报表\par}
{\small\color{refgray}数据：三菱电机此前收购Nozomi Networks获得工厂全域数据入口\par}
{\small\color{refgray}边际变化：三菱电机FA数字化战略从数据采集（Nozomi Networks）延伸至视觉感知层，边缘AI+传感器技术降低数据传输延迟和云端依赖。\par}
\paragraph{GS} LG Display 2Q26营收5.6万亿韩元，营业亏损1080亿韩元（剔除一次性自愿退休成本2400亿韩元后正常化营业利润约1300亿韩元）。3Q ASP指引仅环比增长高十位数百分比，远低于此前预期的30\%环比增幅，主因存储芯片涨价挤压客户预算。智能手机OLED出货量上调至8600万片（同比+13\%），但ASP同比下滑5-8\%。大尺寸OLED受益于世界杯和电竞显示器需求，出货量预计750万片（同比+17\%），广州OLED TV厂折旧结束，大尺寸OLED营业利润率同比改善至双位数。
{\small\color{refgray}数据：2Q26营收5.6万亿韩元，营业亏损1080亿韩元\par}
{\small\color{refgray}数据：3Q ASP指引环比增长高十位数\%，远低于此前预期的30\%\par}
{\small\color{refgray}数据：2026年智能手机OLED出货量预测从8200万片上调至8600万片（同比+13\%）\par}
{\small\color{refgray}数据：智能手机和平板OLED ASP同比下滑5-8\%\par}
{\small\color{refgray}边际变化：出货量上调与ASP下调同时发生，量增价跌组合意味着收入增长可能低于出货量增速；存储芯片涨价成为OLED定价新压制因素。\par}
\subsubsection*{关键数据}
\begin{itemize}
  \item MLCC 2Q26营收442亿新台币，环比+16\%，同比+35\%
  \item 标准无卤板报价>300元/张，2025年底160-170元，2021年高点约280元
  \item 电子布供应偏紧至少持续至2027年底
  \item LG Display 3Q ASP指引环比增长高十位数\%，远低于此前预期的30\%
  \item 智能手机OLED出货量上调至8600万片（同比+13\%），但ASP同比下滑5-8\%
  \item 三菱电机与索尼成立AI视觉传感器合资公司，2026年10月运营
\end{itemize}
\begin{figure}[htbp]
\centering
\includegraphics[width=0.88\linewidth]{figures/fig_086.jpg}
\caption{Exhibit 2 - Exhibit 2: Walsin Tech's preliminary June profits show meaningful margin improvement vs 1Q26 Exhibit 3: Walsin Tech's 1Q26 product mix Product Sales Contributions \& Applications as of Q1 2026 Pre-registration is required. Date an}
\end{figure}
\begin{figure}[htbp]
\centering
\includegraphics[width=0.88\linewidth]{figures/fig_108.jpg}
\caption{Figure 1 - Figure 1. A snapshot of 1H26 preliminary results on CCL / PCB chain Figure 2. E-glass fabric ASP more than double of YTD 2026 © 2026 Citi Inc. No redistribution without Citi's written permission. \#\# Kingboard Laminates Holdings \#\# Valuation Our target price fo}
\end{figure}
\begin{figure}[htbp]
\centering
\includegraphics[width=0.88\linewidth]{figures/fig_139.jpg}
\caption{Exhibit 1 - Exhibit 1: LGD shares are currently trading at 0.7X 12m fwd P/B with 2026E/2027E ROE of -5\%/11\% LGD 12m fwd P/B vs. ROE Exhibit 2: LGD earnings revisions}
\end{figure}
\begin{figure}[htbp]
\centering
\includegraphics[width=0.88\linewidth]{figures/fig_087.jpg}
\caption{Exhibit 4 - Exhibit 4: Quarterly Sales, by Product, 1Q23-1Q26 1Q23 2Q23 3Q23 4Q23 1Q24 2Q24 3Q24 4Q24 1Q25 2Q25 3Q25 4Q25 1Q26 MLCC Chip Resistor Tantalum Magnetics (includes Wireless \& Power) Sensor Other components Exhibit 5: Annual Sales,}
\end{figure}
{\small\color{refgray}本节综合 5 篇研究；涉及 MS、Citi、GS。}
\newpage
\subsection{汽车与出行：盈利分化与结构性转型}
\textbf{汽车行业二季度财报显示，传统车企与新兴科技公司盈利表现分化加剧：特斯拉FSD渗透率加速但毛利率承压，大众受高端品牌残值拖累利润低于预期，Mobileye核心盈利能力仍待验证。结构性转型方向明确，但短期财务兑现存在不确定性。}
\subsubsection*{主流共识}
\begin{itemize}
  \item FSD与自动驾驶技术正在从选项变为购车决策的关键因素，用户使用数据加速增长。
  \item 汽车行业整体面临毛利率压力，受定价、成本及一次性因素影响。
  \item 传统车企与新兴科技公司盈利分化加剧，高端电动车二手价格不确定性影响品牌组利润。
\end{itemize}
\subsubsection*{分歧与条件差异}
\begin{itemize}
  \item 特斯拉毛利率环比持平是否反映定价纪律改善，还是受一次性因素掩盖，市场存在分歧。
  \item Mobileye的研发激励是否可持续，以及其核心业务盈利能力能否改善，观点不一。
  \item 大众下调营收指引但维持利润率目标，市场对其成本削减计划能否抵消销量下滑存疑。
\end{itemize}
\subsubsection*{机构观点}
\paragraph{JPM} 特斯拉二季度调整后EPS 0.33美元低于预期，但FSD订阅用户环比增至150万，渗透率升至55\%，每用户月度行驶里程从390英里加速至540英里。Robotaxi在6城市完成超38万英里无人监督行驶，无事故。毛利率16.3\%低于预期，但剔除一次性因素后环比持平，显示定价纪律未恶化。
{\small\color{refgray}数据：Q2调整后EPS 0.33美元，低于JPM预期0.60美元\par}
{\small\color{refgray}数据：FSD订阅用户从130万增至150万，渗透率55\%\par}
{\small\color{refgray}数据：每用户月度行驶里程从390英里增至540英里\par}
{\small\color{refgray}数据：Robotaxi无人监督行驶超38万英里，无事故\par}
{\small\color{refgray}边际变化：FSD渗透率从30-35\%跳升至55\%，用户行驶里程加速，Robotaxi从测试进入规模化验证阶段。\par}
\paragraph{GS} Mobileye二季度营收5.08亿美元超预期，但调整后毛利率同比下降300bp至65.6\%。剔除以色列政府9300万美元研发激励后，调整后营业利润率仅12.2\%。公司上调全年营收指引至19.7-20.2亿美元，但核心业务指引实则小幅下修。获得Stellantis REM地图订单，Robotaxi计划2027年投放100辆。
{\small\color{refgray}数据：Q2营收5.08亿美元，高于GS预期4.81亿美元\par}
{\small\color{refgray}数据：调整后毛利率65.6\%，同比降300bp\par}
{\small\color{refgray}数据：研发激励贡献9300万美元营业利润\par}
{\small\color{refgray}数据：剔除激励后营业利润率约12.2\%\par}
{\small\color{refgray}边际变化：业绩超预期主要由一次性研发激励驱动，核心业务盈利能力仍待改善；Robotaxi从供应商转向运营方。\par}
\paragraph{GS} 大众二季度运营利润34.7亿欧元，低于市场预期38.5亿欧元，主因Progressive品牌组受残值重估影响利润低于预期29\%。核心品牌组利润超预期8.5\%。汽车业务净现金流11.72亿欧元，较预期5.83亿欧元高101\%，受益于营运资本和资本支出减少。营收增速指引下调至-3\%-0\%，但利润率目标维持4.0-5.5\%。
{\small\color{refgray}数据：Q2运营利润34.7亿欧元，低于预期38.5亿欧元\par}
{\small\color{refgray}数据：Progressive品牌组利润低于预期29\%\par}
{\small\color{refgray}数据：核心品牌组利润超预期8.5\%\par}
{\small\color{refgray}数据：汽车业务净现金流11.72亿欧元，超预期101\%\par}
{\small\color{refgray}边际变化：营收增速指引从0-3\%下调至-3\%-0\%，但利润率目标不变，市场关注成本削减计划能否对冲销量下滑。\par}
\paragraph{HSBC} 福耀玻璃A/H股年初至今下跌15\%/17\%，同期沪深300下跌2\%。国内乘用车批发量前6月同比下降6\%，但福耀盈利韧性较强：国内ASP持续扩张，海外份额提升，成本端纯碱、运费趋于稳定。HSBC下调2026/2027年盈利预测16\%/13\%，但上调毛利率预测0.3/0.9个百分点，新工厂爬坡有望支撑毛利率环比改善。
{\small\color{refgray}数据：福耀A/H股年初至今下跌15\%/17\%\par}
{\small\color{refgray}数据：国内乘用车批发量前6月同比下降6\%\par}
{\small\color{refgray}数据：HSBC下调2026/2027年盈利预测16\%/13\%\par}
{\small\color{refgray}数据：上调毛利率预测0.3/0.9个百分点\par}
{\small\color{refgray}边际变化：盈利预测下调主要反映国内销量变化，但ASP扩张和海外份额增长部分抵消，毛利率预测反而上调。\par}
\subsubsection*{关键数据}
\begin{itemize}
  \item 特斯拉FSD订阅用户环比增至150万，渗透率55\%
  \item Mobileye Q2营收5.08亿美元，超预期5.6\%
  \item 大众Q2运营利润34.7亿欧元，低于预期10\%
  \item 福耀A/H股年初至今下跌15\%/17\%
  \item 国内乘用车批发量前6月同比下降6\%
  \item 特斯拉Robotaxi无人监督行驶超38万英里
  \item Mobileye研发激励贡献9300万美元
  \item 大众汽车业务净现金流11.72亿欧元，超预期101\%
\end{itemize}
\begin{figure}[htbp]
\centering
\includegraphics[width=0.88\linewidth]{figures/fig_043.jpg}
\caption{Figure 1 - Figure 1: TSLA's Vertically and Horizontally Integrated Engineering Playbook Raw materials / upstream inputs Minerals \& raw materials Software / operating loop FSD, vehicle software \& data loop}
\end{figure}
\begin{figure}[htbp]
\centering
\includegraphics[width=0.88\linewidth]{figures/fig_159.jpg}
\caption{Exhibit 1 - Price relative: 3606 HK Price relative: 600660 CH Note: Priced at close of 20 Jul 2026 Key charts Exhibit 1. Resilient GPM and OPM performance during past quarters}
\end{figure}
\begin{figure}[htbp]
\centering
\includegraphics[width=0.88\linewidth]{figures/fig_162.jpg}
\caption{Exhibit 7 - Exhibit 7. Fuyao-A trades at 13.5x 12-month forward P/E vs its historical average of 18.3x Exhibit 8. Fuyao-A trades at 3.3x 12-month forward P/B vs its historical average of 3.1x Exhibit 9. Fuyao-H trades at 11.5x 12-month forward P/E vs its historical averag}
\end{figure}
\begin{figure}[htbp]
\centering
\includegraphics[width=0.88\linewidth]{figures/fig_044.jpg}
\caption{Figure 2 - Figure 2: TSLA's Cross-Platform Shared Hardware Stack}
\end{figure}
{\small\color{refgray}本节综合 4 篇研究；涉及 JPM、GS、HSBC。}
\newpage
\subsection{医疗健康：代谢、手术与AI的价值迁移}
\textbf{GLP-1竞争从减重疗效转向代谢获益与便利性；手术机器人长期渗透空间广阔，短期担忧提供窗口；AI药物发现价值从工具向资产迁移，护城河在于数据闭环而非模型。}
\subsubsection*{主流共识}
\begin{itemize}
  \item GLP-1下一代疗法需实现>20\%安慰剂调整减重，并验证心血管、肾脏等代谢获益。
  \item 口服GLP-1（orforglipron、口服司美格鲁肽）可实现约10\%减重，拓宽治疗人群。
  \item 软组织手术机器人渗透率仍低，新应用场景（心脏、良性普外、日间手术）驱动长期增长。
  \item AI药物发现货币化仍集中在低价值层（SaaS、CRO），长期价值向下游管线资产迁移。
\end{itemize}
\subsubsection*{分歧与条件差异}
\begin{itemize}
  \item GLP-1口服剂型中，小分子（orforglipron）vs 肽类（口服司美格鲁肽）的依从性优势存在分歧。
  \item 手术机器人领域，新进入者（强生Ottava）能否突破Intuitive生态壁垒观点不一。
  \item AI药物发现公司估值框架：市场定价由期权价值驱动，但传统DCF是否适用存争议。
\end{itemize}
\subsubsection*{机构观点}
\paragraph{Bernstein} GLP-1市场进入代谢获益与便利性竞争新阶段。下一代注射疗法（retatrutide、amycretin等）将减重基准推至>20\%，口服GLP-1（orforglipron、口服司美格鲁肽）实现约10\%减重。小分子口服药（orforglipron）无需空腹，依从性优势显著。信达IBI3032四周减重8.6\%，周制剂口服候选IBI3042体现降频趋势。
{\small\color{refgray}数据：retatrutide第80周安慰剂调整减重约26.1\%\par}
{\small\color{refgray}数据：皮下amycretin第36-48周减重约23.2\%\par}
{\small\color{refgray}数据：orforglipron和口服司美格鲁肽减重约10\%\par}
{\small\color{refgray}数据：信达IBI3032四周减重8.6\%\par}
{\small\color{refgray}边际变化：竞争焦点从单纯减重疗效扩展至代谢获益、肌肉保留和给药便利性。\par}
\paragraph{Bernstein} 软组织手术机器人长期渗透空间广阔，短期担忧提供窗口。AI正提高进入壁垒，Intuitive构建的生态系统（集成系统、软件、数据、培训）使新进入者需解决未满足问题才能获得 traction。强生Ottava获批，为首个与手术台集成平台，占地面积缩小30\%-50\%，但需证明差异化价值。
{\small\color{refgray}数据：Intuitive投入数十亿美元构建生态系统\par}
{\small\color{refgray}数据：强生Ottava占地面积缩小30\%-50\%\par}
{\small\color{refgray}数据：新应用场景包括心脏、非工作时间、良性普外、日间手术\par}
{\small\color{refgray}边际变化：市场关注点从装机量增速转向长期渗透率提升，AI成为新竞争壁垒。\par}
\paragraph{Citi} 微创医疗机器人首次实现半年度正向利润，营收同比增长200\%-230\%，超预期。海外营收同比增长超450\%，成为增长主引擎。毛利率同比提升超15个百分点，得益于规模效应和高毛利耗材迁移。利润中值3400万元低于预期4400万元，归因于研究驱动型增长阶段。
{\small\color{refgray}数据：净利润2800-4000万元，去年同期亏损1.15亿元\par}
{\small\color{refgray}数据：营收同比增长200\%-230\%\par}
{\small\color{refgray}数据：海外营收同比增长超450\%\par}
{\small\color{refgray}数据：毛利率同比提升超15个百分点\par}
{\small\color{refgray}边际变化：首次半年度盈利，海外营收增速远超国内，管理层可能上调全年装机目标。\par}
\paragraph{GS} ESMO 2026中国数据聚焦：科伦博泰sac-TMT在PD-L1 TPS<1\%一线非鳞NSCLC中达到PFS主要终点，若入选LBA将是焦点。再鼎医药ZL-1310一线SCLC早期数据、多个KRAS G12D抑制剂胰腺癌首次人体数据值得关注。ADC联合治疗的安全性与疗效平衡是关键观察点。
{\small\color{refgray}数据：sac-TMT在PD-L1 TPS<1\%亚群达到PFS主要终点\par}
{\small\color{refgray}数据：LBA摘要标题9月25日公布，数据10月19日披露\par}
{\small\color{refgray}数据：ZL-1310为DLL3靶向疗法一线SCLC\par}
{\small\color{refgray}边际变化：中国创新药在ADC联合治疗和RAS靶向领域进入数据验证密集期。\par}
\paragraph{GS} 罗氏/高德美二季度业绩为日本医药提供参照。Hemlibra全球增长持续，优于市场共识4\%；Nemluvio表现高于预期。罗氏重申中国医疗定价改革影响2026年趋于平稳，与雅培、丹纳赫观点一致，利好希森美康。GYM329肥胖试验数据读出按计划2026年进行。
{\small\color{refgray}数据：Hemlibra优于市场共识4\%\par}
{\small\color{refgray}数据：Nemluvio表现高于预期\par}
{\small\color{refgray}数据：中国医疗定价改革影响预计2026年平稳\par}
{\small\color{refgray}数据：GYM329肥胖试验数据读出2026年\par}
{\small\color{refgray}边际变化：中国诊断定价压力缓解信号强化，肥胖管线进展按计划推进。\par}
\paragraph{HSBC} AI药物发现价值从工具向资产迁移。当前货币化集中在低价值层（SaaS、CRO），长期价值向下游管线所有权迁移。护城河不在模型，而在专有数据生成、工作流整合、湿实验室基础设施形成的反馈闭环。市场定价由期权价值驱动，传统DCF不适用。
{\small\color{refgray}数据：货币化集中在低到中等价值层\par}
{\small\color{refgray}数据：护城河来自反馈闭环而非算法\par}
{\small\color{refgray}数据：纯SaaS模式面临商品化风险\par}
{\small\color{refgray}边际变化：市场认知从AI工具公司转向平台型生物科技，估值框架从DCF转向期权定价。\par}
\subsubsection*{关键数据}
\begin{itemize}
  \item retatrutide第80周安慰剂调整减重约26.1\%
  \item orforglipron和口服司美格鲁肽减重约10\%
  \item 微创机器人海外营收同比增长超450\%
  \item sac-TMT在PD-L1 TPS<1\% NSCLC达到PFS主要终点
  \item Hemlibra优于市场共识4\%
  \item AI药物发现货币化集中在低价值层
\end{itemize}
\begin{figure}[htbp]
\centering
\includegraphics[width=0.88\linewidth]{figures/fig_006.jpg}
\caption{EXHIBIT 1 - EXHIBIT 1: Weight loss, convenience, body composition/metabolic benefits are the key investor preference when choosing GLP-1 products \#\# PART 1: WEIGHT LOSS EFFICACY AND EFFICIENCY Efficacy benchmark in weight loss has evolved fr}
\end{figure}
\begin{figure}[htbp]
\centering
\includegraphics[width=0.88\linewidth]{figures/fig_154.jpg}
\caption{Exhibit 2 - 4. What becomes the long-term moat? We expect most AI advantages to become commoditised. Instead, durable moats will emerge from proprietary data generation, workflow integration, regulatory positioning, and wet-lab infrastructure – not simply better algorithm}
\end{figure}
\begin{figure}[htbp]
\centering
\includegraphics[width=0.88\linewidth]{figures/fig_008.jpg}
\caption{Exhibit 3 - EXHIBIT 4: Large overlapping population between obesity, diabetes, CKD, CVD, and MASH}
\end{figure}
\begin{figure}[htbp]
\centering
\includegraphics[width=0.88\linewidth]{figures/fig_156.jpg}
\caption{Exhibit 6 - The moat is not the model. The moat is the Feedback Loop. \#\# Exhibit 6. Global AIDD market trend: US and China are building ecosystems with different approaches Global AIDD deal trends China spotlight: Value migration Frontier AI \& big tech: The biology vertic}
\end{figure}
{\small\color{refgray}本节综合 6 篇研究；涉及 Bernstein、Citi、GS、HSBC。}
\newpage
\subsection{中国消费与零售：结构性分化与品牌再分配}
\textbf{中国消费市场整体增速放缓，但结构性机会显著：本土美妆品牌份额持续扩大但内部分化加剧，黄金珠宝受金价波动驱动增长但固定价格产品承压，德系豪华车量跌价分化，音乐流媒体高频使用但付费转化仍是关键。}
\subsubsection*{主流共识}
\begin{itemize}
  \item 中国美妆市场增长动力从总量扩张转向品牌间份额再分配，本土品牌在护肤品类市占率持续提升，但增速分化明显。
  \item 黄金珠宝行业同店增长由金价回调驱动的重量计价产品带动，固定价格产品仍承压，消费者更注重黄金保值属性。
  \item 德系豪华车在华零售量普遍两位数下降，但进口车型价格韧性优于国产，宝马M系列实现量价齐升。
\end{itemize}
\subsubsection*{分歧与条件差异}
\begin{itemize}
  \item 美妆品牌内部：Proya等头部品牌在电商大促中持续领先，而Giant Biogene等增速骤降，Bloomage持续下滑，显示品牌分层加速。
  \item 黄金珠宝：周大福固定价格产品占比同比下降约5个百分点，与全年指引中小幅扩张的目标存在差距，新产品能否扭转是关键分歧。
  \item 音乐流媒体：Soda Music以免费模式快速获取用户（渗透率43\%），但月均消费12.1元低于行业平均14-15元，付费转化路径存疑。
\end{itemize}
\subsubsection*{机构观点}
\paragraph{MS} 中国美妆市场进入存量竞争阶段，本土品牌在护肤品类市占率从2012年的约21\%升至2024年超50\%，但内部分化显著。Proya在2024-2025年多个季度保持正向增长，而Giant Biogene增速从2024H1的58\%骤降至2025H2的-20\%，Bloomage功能性护肤业务持续下滑。国际品牌在高端线和抖音渠道仍具竞争力，进口品牌正调整产品组合与定价策略以适应市场变化。
{\small\color{refgray}数据：本土品牌护肤市占率从2012年约21\%升至2024年超50\%\par}
{\small\color{refgray}数据：Proya季度销售额同比增速从2021年30\%-49\%降至2025年-12\%至8\%\par}
{\small\color{refgray}数据：Giant Biogene 2025H2增速骤降至-20\%\par}
{\small\color{refgray}数据：Bloomage功能性护肤业务2022年起持续下滑\par}
{\small\color{refgray}边际变化：此前市场关注总量增长，现转向品牌间份额再分配与内部分化。\par}
\paragraph{Citi} 中国音乐流媒体用户日均使用88分钟，66\%日活跃率接近短视频的70\%，18-29岁用户日均118分钟。Soda Music渗透率达43\%，但41\%用户仅用免费版，月均消费12.1元低于腾讯音乐和网易云音乐的14-15元。用户选择Soda Music首要原因是免费（32\%），曲库规模与音质是切换平台主因。AI音乐接受度提升，70\%受访者听过，43\%难以区分AI与原创。
{\small\color{refgray}数据：用户日均使用88分钟，66\%日活跃率\par}
{\small\color{refgray}数据：18-29岁用户日均118分钟\par}
{\small\color{refgray}数据：Soda Music渗透率43\%，41\%仅用免费版\par}
{\small\color{refgray}数据：Soda Music月均消费12.1元 vs 行业14-15元\par}
{\small\color{refgray}边际变化：新进入者以免费模式快速获取用户，但付费转化和ARPU提升成为关键关注点。\par}
\paragraph{GS} 周大福1QFY27中国内地自营店同店销售增长19.6\%，港澳增长41.7\%，主要由金价回调驱动的重量计价产品带动（内地+38\%，港澳+64\%）。固定价格产品内地同比下降1\%，占比同比降约5个百分点。7月前两周受天气和世界杯分流影响，增速回落至个位数，第三周已环比改善。Dawn系列首季零售额超8.7亿港元，新客占比超20\%，全年销售预期上调至30亿港元。
{\small\color{refgray}数据：中国内地自营店同店销售增长19.6\%\par}
{\small\color{refgray}数据：港澳同店销售增长41.7\%\par}
{\small\color{refgray}数据：重量计价产品内地+38\%，港澳+64\%\par}
{\small\color{refgray}数据：固定价格产品内地同比-1\%，占比降约5个百分点\par}
{\small\color{refgray}边际变化：此前市场关注金价回调对量的提振，现转向固定价格产品占比能否回升及7月增速放缓后的可持续性。\par}
\paragraph{GS} 德系豪华车6月零售量普遍两位数下降：奔驰国产-41.8\%，进口-17.9\%；宝马国产-32.7\%，进口-45.6\%；保时捷-46.2\%。但价格端分化：宝马M系列销量+9.9\%，均价+5.4\%至82.9万元，实现量价齐升；奔驰进口均价降至90.8万元（-11.6\%），S级销量1,059辆（-38.8\%），迈巴赫版成交价下降。本土竞品尊界S800从4,000+峰值回落至593辆，波动性显著。
{\small\color{refgray}数据：奔驰国产-41.8\%，进口-17.9\%\par}
{\small\color{refgray}数据：宝马国产-32.7\%，进口-45.6\%\par}
{\small\color{refgray}数据：保时捷-46.2\%\par}
{\small\color{refgray}数据：宝马M系列销量+9.9\%，均价+5.4\%至82.9万元\par}
{\small\color{refgray}边际变化：此前市场关注整体量价下行，现发现宝马M系列逆势量价齐升，本土高端电动车销量波动性远高于传统豪华燃油车。\par}
\subsubsection*{关键数据}
\begin{itemize}
  \item 本土品牌护肤市占率从2012年约21\%升至2024年超50\%
  \item 周大福中国内地自营店同店销售增长19.6\%
  \item 音乐流媒体用户日均使用88分钟，66\%日活跃率
  \item Soda Music渗透率43\%，月均消费12.1元低于行业14-15元
  \item 宝马M系列销量+9.9\%，均价+5.4\%至82.9万元
  \item 奔驰国产-41.8\%，进口-17.9\%
  \item Dawn系列首季零售额超8.7亿港元，全年预期上调至30亿
  \item 70\%受访者听过AI音乐，43\%难以区分AI与原创
\end{itemize}
\begin{figure}[htbp]
\centering
\includegraphics[width=0.88\linewidth]{figures/fig_114.jpg}
\caption{Exhibit 1 - Exhibit 1: Chow Tai Fook Jun-quarter SSSG accelerated in ML China, and remained strong in HK \& Macau Chow Tai Fook SSSG by region Note: Normal period for HK\&Macau market denotes pre-COVID and before HK macro situation \#\# Price Ta}
\end{figure}
\begin{figure}[htbp]
\centering
\includegraphics[width=0.88\linewidth]{figures/fig_115.jpg}
\caption{Exhibit 1 - Exhibit 1: Mercedes JV volume: 24,800 units sold in Jun-26, R1/3/6/12M at -41.8\%/-32.8\%/-31.2\%/-28.3\% yoy. Exhibit 2: Mercedes import volume: 6,869 units sold in Jun-26 (R1/3/6/12M -17.9\%/-22.3\%/-24.1\%/-25.0\%), of which TEV a}
\end{figure}
\begin{figure}[htbp]
\centering
\includegraphics[width=0.88\linewidth]{figures/fig_101.jpg}
\caption{Figure 1 - ■ Demographic — 15\% of respondents are aged 12–17, 25\% are aged 18–29, 25\% are aged 30–39, and 35\% are aged 40 and above. The mean age of the sample is 33.5 years. ■ Geographic distribution — Among the respondents, 33\% reside in Tier-1 cities, 33\% in Tier-2 ci}
\end{figure}
\begin{figure}[htbp]
\centering
\includegraphics[width=0.88\linewidth]{figures/fig_102.jpg}
\caption{Figure 1 - ■ Content consumption profile — In addition to music streaming, 86\% also consume short-video content, 79\% stream movies/shows, 65\% engage with livestreaming, 61\% play video games via streaming, 57\% listen to audiobooks, and 44\% consume podcasts via streaming p}
\end{figure}
{\small\color{refgray}本节综合 4 篇研究；涉及 MS、Citi、GS。}
\newpage
\subsection{能源与公用事业：结构性转型与地缘扰动下的再平衡}
\textbf{日本电力需求因AI和半导体迎来结构性增长，但发电转型节奏决定定价逻辑；油轮运价从高位回落至均衡区，地缘事件和需求边际变化成为关键变量；中国多晶硅能耗标准收紧更多影响闲置产能，短期供需宽松格局不变。}
\subsubsection*{主流共识}
\begin{itemize}
  \item 日本电力需求正经历结构性变化，AI数据中心和半导体工厂是主要增量来源，预计2041年发电量增至1.1-1.2万亿千瓦时。
  \item 油轮即期运价已从6月高点大幅回落，大西洋市场VLCC运价接近均衡水平，进一步下行空间有限。
  \item 中国多晶硅行业产能利用率仅30-40\%，新能耗标准主要淘汰已闲置的低效产能，对有效供给影响微乎其微。
\end{itemize}
\subsubsection*{分歧与条件差异}
\begin{itemize}
  \item 日本发电结构转型节奏存在不确定性：可再生能源目标40-50\%与核能20\%的达成路径，取决于政策落地速度和技术进展。
  \item 油轮市场后续走势分歧在于：地缘事件（中东遇袭、红海绕行）持续消耗有效运力，但中国进口量处于低位，需求回暖时机不明。
  \item 中国多晶硅能耗标准收紧是否预示更严格的供给调控：当前观点认为仅影响闲置产能，但若未来执行力度加强，可能改变行业格局。
\end{itemize}
\subsubsection*{机构观点}
\paragraph{MS} 日本电力需求因AI数据中心和半导体工厂新建而结构性增长，第七次能源基本计划预计2041年发电量达1.1-1.2万亿千瓦时，较2024年9854亿千瓦时显著提升。发电结构将大幅调整：可再生能源占比从23\%升至40-50\%，核能从9\%升至20\%，火电从69\%压缩至30-40\%。但核电重启进度不一，西部快于东部，仅15座机组恢复运行。能源自给率目标从15.2\%提升至约30\%。
{\small\color{refgray}数据：日本2041年发电量目标1.1-1.2万亿千瓦时，2024年为9854亿千瓦时\par}
{\small\color{refgray}数据：可再生能源目标占比40-50\%，核能20\%，火电30-40\%\par}
{\small\color{refgray}数据：33座核电机组中仅15座恢复运行\par}
{\small\color{refgray}数据：能源自给率目标从15.2\%提升至约30\%\par}
{\small\color{refgray}边际变化：日本电力需求从过去十年增长停滞转向结构性增长，核心驱动力从传统工业转向AI和半导体政策优先领域。\par}
\paragraph{MS} 油轮即期运价从6月下旬高点快速回落后趋于均衡，7月21日西非-中国航线TD15报10.3万美元/天，美湾-中国航线TD22报10.5万美元/天，较6月底18.9万和15.5万显著回落。地缘事件（中东油轮遇袭、黑海CPC码头受损）及红海绕行消耗有效运力，而中国6月原油进口2927万吨（同比-41.3\%）为2016年10月以来新低，需求端是最大变量。
{\small\color{refgray}数据：7月21日TD15运价10.3万美元/天，TD22运价10.5万美元/天\par}
{\small\color{refgray}数据：6月下旬TD15运价18.9万美元/天，TD22运价15.5万美元/天\par}
{\small\color{refgray}数据：中国6月原油进口2927万吨，同比-41.3\%，为2016年10月以来新低\par}
{\small\color{refgray}边际变化：油轮市场从‘低风险高回报’转向‘回报与风险更均衡’状态，地缘事件持续扰动供给，但需求端中国进口疲软成为关键制约。\par}
\paragraph{Citi} 7月21日多晶硅能耗限额国标正式版发布，棒状硅三级能耗上限从6.4降至6.3kgce/kg，颗粒硅从5.0降至4.6kgce/kg。前四大多晶硅供应商（通威、协鑫、大全、特变/新特能源）合计名义产能191万吨均能达标，行业平均产能利用率仅30-40\%，新标准主要淘汰已闲置的低效产能。7月多晶硅月产量预计10万吨，月需求9.1万吨，供给过剩约10\%。
{\small\color{refgray}数据：棒状硅三级能耗上限从6.4降至6.3kgce/kg，颗粒硅从5.0降至4.6kgce/kg\par}
{\small\color{refgray}数据：前四大多晶硅供应商合计名义产能191万吨，均能达标\par}
{\small\color{refgray}数据：行业平均产能利用率30-40\%\par}
{\small\color{refgray}数据：7月多晶硅月产量10万吨，月需求9.1万吨，供给过剩约10\%\par}
{\small\color{refgray}边际变化：能耗标准收紧方向明确，但主要影响已闲置产能，对当前有效供给影响微乎其微，市场不应将其解读为供给收缩信号。\par}
\subsubsection*{关键数据}
\begin{itemize}
  \item 日本2041年发电量目标1.1-1.2万亿千瓦时，2024年为9854亿千瓦时
  \item 日本可再生能源目标占比40-50\%，核能20\%，火电30-40\%
  \item 7月21日TD15运价10.3万美元/天，较6月下旬18.9万回落
  \item 中国6月原油进口2927万吨，同比-41.3\%，为2016年10月以来新低
  \item 棒状硅三级能耗上限从6.4降至6.3kgce/kg，颗粒硅从5.0降至4.6kgce/kg
  \item 前四大多晶硅供应商合计名义产能191万吨，均能达标
  \item 行业平均产能利用率30-40\%
  \item 7月多晶硅月产量10万吨，月需求9.1万吨，供给过剩约10\%
\end{itemize}
\begin{figure}[htbp]
\centering
\includegraphics[width=0.88\linewidth]{figures/fig_083.jpg}
\caption{Exhibit 1 - Exhibit 1: Tanker spot earnings have corrected again since mid-June Exhibit 2: China crude imports have been soft since April 2026 Qianlei Fan, CFA}
\end{figure}
\begin{figure}[htbp]
\centering
\includegraphics[width=0.88\linewidth]{figures/fig_106.jpg}
\caption{Figure 1 - © 2026 Citi Inc. No redistribution without Citi's written permission. Figure 1. PRC weekly polysilicon price © 2026 Citi Inc. No redistribution without Citi's written permission. Figure 2. Monthly polysilicon output in China \#\# Companies Mentioned:}
\end{figure}
\begin{figure}[htbp]
\centering
\includegraphics[width=0.88\linewidth]{figures/fig_084.jpg}
\caption{Exhibit 1 - Exhibit 1: Tanker spot earnings have corrected again since mid-June Exhibit 2: China crude imports have been soft since April 2026 Qianlei Fan, CFA}
\end{figure}
\begin{figure}[htbp]
\centering
\includegraphics[width=0.88\linewidth]{figures/fig_107.jpg}
\caption{Figure 1 - Figure 1. PRC weekly polysilicon price © 2026 Citi Inc. No redistribution without Citi's written permission. Figure 2. Monthly polysilicon output in China \#\# Companies Mentioned: Daqo New Energy (DQ.N; US\textbackslash{}\$11.86; 1H; 21 Jul 26; 16:00) GCL Technology (3800.HK; }
\end{figure}
{\small\color{refgray}本节综合 3 篇研究；涉及 MS、Citi。}
\newpage
\subsection{工业与材料：数据中心供需紧张与智驾渗透加速}
\textbf{数据中心受益于AI需求的结构性窗口，续租价差与积压订单显示定价权强劲，但股权稀释构成短期压力；智能驾驶领域，地平线HSD 2.0推动高阶智驾下沉至10-20万元车型，市场份额快速提升，与大众合作向L3/L4演进。}
\subsubsection*{主流共识}
\begin{itemize}
  \item 数据中心行业供需紧张，大容量客户续租溢价显著，反映优质空间稀缺。
  \item AI应用推动数据中心需求持续增长，积压订单提供中期可见性。
  \item 智能驾驶方案向中低端车型渗透，地平线等供应商市场份额快速提升。
\end{itemize}
\subsubsection*{分歧与条件差异}
\begin{itemize}
  \item 对数据中心股权稀释影响的评估：Bernstein认为1200万股增发几乎抵消2026年每股回报增长，而GS更关注运营指标超预期。
  \item 对地平线净利润波动的解读：可转换融资公允价值变动导致调整后净利润亏损，但经营层面增长强劲，市场对此分歧。
  \item 对数据中心企业托管业务增长可持续性的看法：Bernstein强调其定价能力和利润率优势，但GS未单独突出此业务。
\end{itemize}
\subsubsection*{机构观点}
\paragraph{Bernstein} Digital Realty二季度核心运营指标超预期约6\%，但7月1日与黑石交易相关的1200万股增发几乎完全抵消2026年每股回报增长。1.88亿美元净促进收益暂时缓解冲击，但不可持续。企业托管业务首次突破1亿美元季度收入，互连服务贡献19\%，客户流失率低，被视为增长轨迹中最值得关注的部分。超大规模续租现金价差66\%、GAAP价差92\%，定价能力强劲。在建项目1.4GW，54\%已预租，全球可建容量超8.5GW。
{\small\color{refgray}数据：二季度核心运营指标超预期约6\%\par}
{\small\color{refgray}数据：7月1日发行1200万股，几乎抵消2026年每股回报增长\par}
{\small\color{refgray}数据：企业托管业务季度收入首次突破1亿美元\par}
{\small\color{refgray}数据：互连服务占企业托管收入约19\%\par}
{\small\color{refgray}边际变化：此前市场关注运营增长，Bernstein新增对股权稀释的量化分析，指出增发对每股回报的抵消效应。\par}
\paragraph{Bernstein} 霍尼韦尔2026财年二季度有机收入同比增长4\%，高于市场预期的持平；调整后每股盈利1.95美元，比共识高7\%。三大业务板块均优于预期：过程自动化有机增长-1\%，远好于预期的-7.5\%；楼宇自动化有机增长9\%，高于预期的6.5\%；工业自动化有机增长4\%，高于预期的0.5\%。分部利润率均高于预期0.2-0.3个百分点。全年指引小幅上调：有机增长预期从2-3\%升至3-4\%，每股盈利中值从8.10升至8.20美元。
{\small\color{refgray}数据：有机收入同比增长4\%，高于预期的持平\par}
{\small\color{refgray}数据：调整后每股盈利1.95美元，比共识高7\%\par}
{\small\color{refgray}数据：过程自动化有机增长-1\%，优于预期的-7.5\%\par}
{\small\color{refgray}数据：楼宇自动化有机增长9\%，高于预期的6.5\%\par}
{\small\color{refgray}边际变化：此前市场对过程自动化板块预期悲观，实际下滑速度远低于预期，推动全年指引上调。\par}
\paragraph{GS} Digital Realty二季度核心FFO每股2.65美元，高于共识预期的1.98美元；剔除一次性收益后为2.13美元，仍超预期。租赁收入11.46亿美元，同比增长19.3\%，超预期4.7\%。续租年化租金2.61亿美元，现金续租价差25\%，其中0-1MW价差仅5\%，而大于1MW价差高达67\%，反映新加坡等核心市场供需紧张。新订单2.08亿美元，7月再签2.05亿美元超大规模租约。2026年核心FFO指引从8.00-8.10上调至8.15-8.20美元。
{\small\color{refgray}数据：二季度核心FFO每股2.65美元，共识预期1.98美元\par}
{\small\color{refgray}数据：租赁收入11.46亿美元，同比增长19.3\%\par}
{\small\color{refgray}数据：现金续租价差25\%，其中>1MW价差67\%\par}
{\small\color{refgray}数据：新订单2.08亿美元，7月再签2.05亿美元\par}
{\small\color{refgray}边际变化：此前市场关注订单量，GS强调续租价差的结构性分化，67\%的大容量价差凸显供需紧张。\par}
\paragraph{GS} 地平线1H26营收指引中位数200亿元，同比增长25-35\%，同期中国新能源汽车出货量同比下降14\%，增速为行业两倍以上。HSD V2.0基于Journey 6P平台和升级的世界模型与端到端模型，新增18项功能，已搭载于奇瑞iCAR V27，用户采用率高，帮助10-20万元车型实现城市NOA。中国ADAS与自动驾驶市场价值份额从2022年2\%升至2025年9\%，预计2026年达16\%，2028年达37\%。与大众合作向L3/L4迁移，成为2H26及2027年增长驱动。
{\small\color{refgray}数据：1H26营收指引中位数200亿元，同比增长25-35\%\par}
{\small\color{refgray}数据：同期中国新能源汽车出货量同比下降14\%\par}
{\small\color{refgray}数据：HSD V2.0新增18项功能，搭载于奇瑞iCAR V27\par}
{\small\color{refgray}数据：中国ADAS与自动驾驶市场价值份额2022年2\%，2025年9\%，预计2028年37\%\par}
{\small\color{refgray}边际变化：此前市场关注HSD V2.0技术升级，GS新增对市场份额跃升的量化预测，并强调与大众合作向高阶自动驾驶演进。\par}
\subsubsection*{关键数据}
\begin{itemize}
  \item Digital Realty超大规模续租现金价差66\%，GAAP价差92\%
  \item Digital Realty企业托管业务季度收入首次突破1亿美元
  \item Digital Realty二季度核心FFO每股2.65美元，共识预期1.98美元
  \item Digital Realty>1MW续租价差67\%，0-1MW仅5\%
  \item 霍尼韦尔有机收入同比增长4\%，高于预期的持平
  \item 霍尼韦尔过程自动化有机增长-1\%，优于预期的-7.5\%
  \item 地平线1H26营收指引中位数200亿元，同比增长25-35\%
  \item 地平线中国ADAS与自动驾驶市场价值份额预计2028年达37\%
\end{itemize}
\begin{figure}[htbp]
\centering
\includegraphics[width=0.88\linewidth]{figures/fig_011.jpg}
\caption{EXHIBIT 1 - EXHIBIT 3: Impact of One-Time Adjustments and BX Dilution on Adj. Diluted EPS}
\end{figure}
\begin{figure}[htbp]
\centering
\includegraphics[width=0.88\linewidth]{figures/fig_013.jpg}
\caption{EXHIBIT 2 - EXHIBIT 2: Impact of One-Time Adjustments and BX Dilution on Core FFO per Share (Excluding Net Promote) Impact of One-Time Adjustments and BX Dilution- Core FFOPS (excl. Promote) EXHIBIT 3: Impact of One-Time Adjustments and BX D}
\end{figure}
\begin{figure}[htbp]
\centering
\includegraphics[width=0.88\linewidth]{figures/fig_135.jpg}
\caption{Exhibit 3 - Exhibit 3: The target multiple is derived from peers EV/EBITDA and EBITDA growth Exhibit 4: Horizon Robotics 12M forward P/S}
\end{figure}
\begin{figure}[htbp]
\centering
\includegraphics[width=0.88\linewidth]{figures/fig_012.jpg}
\caption{EXHIBIT 1 - EXHIBIT 3: Impact of One-Time Adjustments and BX Dilution on Adj. Diluted EPS Impact of One-Time Adjustments and BX Dilution- Adj. Diluted EPS \#\# THE LONG VIEW}
\end{figure}
{\small\color{refgray}本节综合 4 篇研究；涉及 Bernstein、GS。}
\newpage
\subsection{区域市场与主题：政策分化与结构性机会}
\textbf{全球市场正经历政策预期与结构性调整的分化：中国财政节奏偏缓但留有后手，日本央行静待10月加息窗口，欧洲央行9月加息概率上升，而美国债市已提前定价利率风险。同时，AI硬件出口、液冷散热、韩国造船等结构性主题提供独立于宏观的alpha机会。}
\subsubsection*{主流共识}
\begin{itemize}
  \item 中国财政政策下半年将加速落地，但短期不会推出类似2024年9月的大规模刺激。
  \item 日本央行7月按兵不动，市场普遍预期10月为下一次加息窗口。
  \item 欧洲央行9月加息概率上升，能源价格是关键变量。
\end{itemize}
\subsubsection*{分歧与条件差异}
\begin{itemize}
  \item 中国自由流动性指标创阶段新低，但MS认为年底前将逐步改善，而改善幅度受企业借贷意愿制约。
  \item 市场对苹果九月季度毛利率预期（47.4\%）与MS估算（46.8\%）存在分歧，服务收入增速亦低于共识。
  \item 韩国造船板块估值（0.5-0.6x P/B）已计入周期中段预期，但海军/AI数据中心引擎等期权价值是否被低估存在分歧。
\end{itemize}
\subsubsection*{机构观点}
\paragraph{BofA} 美国30年期国债实际回报升至3\%（2008年11月以来最高），FCI收紧幅度已超过EPS增长，市场主要矛盾从盈利转向资金成本。BofA牛熊指标维持在9.6的极高水平，触发谨慎信号。83\%的投资者认为美联储11月前不会加息，但市场隐含7月加息概率已升至38\%，9月加息已被完全定价，预期差是资产重新定价的催化剂。
{\small\color{refgray}数据：30年期国债回报5.2\%，创2007年6月以来新高\par}
{\small\color{refgray}数据：30年期实际回报3\%，2008年11月以来最高\par}
{\small\color{refgray}数据：BofA牛熊指标9.6\par}
{\small\color{refgray}数据：7月加息概率38\%\par}
{\small\color{refgray}边际变化：市场从关注盈利转向关注资金成本，FCI收紧速度超过EPS增长，债券市场已提前定价Warsh式政策调整。\par}
\paragraph{HSBC} 全球对中国贸易的推回力度2019-2025年增长75\%，但中国仍是80多个地区的最大进口来源国。贸易流向东移：对美出口占比下降7个百分点，对东盟上升5个百分点。中国在关键矿产和能源转型产品上仍具不可替代性，西方若完全脱钩需投入约24万亿美元。贸易顺差结构分化：对美顺差收窄，对欧顺差扩大。
{\small\color{refgray}数据：贸易限制措施增加75\%\par}
{\small\color{refgray}数据：中国对美出口占比下降7个百分点\par}
{\small\color{refgray}数据：对东盟出口占比上升5个百分点\par}
{\small\color{refgray}数据：西方脱钩成本约24万亿美元\par}
{\small\color{refgray}边际变化：贸易摩擦从关税转向非关税壁垒，但中国出口目的地再平衡至东盟和欧盟，供应链轴心从跨太平洋转向亚洲内部。\par}
\paragraph{JPM} 日本央行7月会议预计按兵不动，但可能上调FY26 GDP增长预测（4月从1.0\%下调至0.5\%）。产出缺口自2022年初持续为正，若增长和通胀预测高于潜在水平，央行需主动回应物价稳定关切。JPM预计10月为下一次加息窗口。
{\small\color{refgray}数据：FY26 GDP增长预测4月从1.0\%下调至0.5\%\par}
{\small\color{refgray}数据：产出缺口自2022年初持续为正\par}
{\small\color{refgray}数据：JPM预计10月加息\par}
{\small\color{refgray}边际变化：日本央行可能提前修正增长预测，以缓解市场对其滞后于曲线的看法，加息节奏取决于如何定义确认每次加息效果所需时间。\par}
\paragraph{JPM} 中国6月一般公共预算收入同比增8.7\%（2025年1月以来最快），税收改善（企业所得税+29.7\%，印花税+145.3\%）。但政府性基金收入同比降31.1\%，土地出让收入降42.1\%（十多年来最大单月降幅）。上半年财政存款增量仍高于近年均值，专项债发行偏慢，下半年政策空间充裕，短期重点在加速预算落地而非增量刺激。
{\small\color{refgray}数据：6月一般公共预算收入同比+8.7\%\par}
{\small\color{refgray}数据：企业所得税+29.7\%\par}
{\small\color{refgray}数据：印花税+145.3\%\par}
{\small\color{refgray}数据：政府性基金收入-31.1\%\par}
{\small\color{refgray}边际变化：税收改善与PPI回升相关，但基建支出持续负增长，资金向项目端传导仍存堵点。二季度执行偏缓为下半年留出政策空间。\par}
\paragraph{JPM} 亚洲组件（AVC）Q2毛利率预期从31\%上调至32.5\%，远高于市场一致预期的约30\%，良率提升和生产效率改善是主因。VR200冷板出货延迟至8月，营收动能有望恢复。Q3营收预计环比+15\%，Q4 AWS Trainium 3和谷歌TPU v8液冷版本量产，单机价值量更高。上调2026/2027 EPS预期4\%和10\%。
{\small\color{refgray}数据：Q2毛利率预期32.5\% vs 市场一致预期约30\%\par}
{\small\color{refgray}数据：Q2 EPS预期24.5新台币 vs 市场一致预期21.4新台币\par}
{\small\color{refgray}数据：Q3营收环比+15\%\par}
{\small\color{refgray}数据：Q4营收环比+16\%\par}
{\small\color{refgray}边际变化：市场低估了AVC的利润弹性，毛利率提升来自良率控制而非价格因素，液冷项目从Q3起进入量产爬坡。\par}
\paragraph{JPM} 韩国造船板块在加拿大潜艇订单落选后调整，韩华海洋-34\%，HD现代重工-31\%，三星重工-21\%，KSOE-15\%。当前P/B在0.5-0.6x，历史峰值约1.0x，市场已计入周期中段预期但未给海军/AI数据中心引擎等期权价值溢价。Q2业绩焦点：商业船回报率能否稳定、引擎产能瓶颈是否缓解、期权价值能否兑现。
{\small\color{refgray}数据：韩华海洋月跌34\%\par}
{\small\color{refgray}数据：HD现代重工月跌31\%\par}
{\small\color{refgray}数据：三星重工月跌21\%\par}
{\small\color{refgray}数据：KSOE月跌15\%\par}
{\small\color{refgray}边际变化：市场关注点从短期表现转向期权价值兑现路径，二季度业绩沟通会能否提供海军/AI数据中心引擎等潜在价值的兑现路径是关键。\par}
\paragraph{MS} Alphabet云业务Q2同比增82\%（超MS预期77\%），EBIT利润率36\%（超预期35\%），增量利润率54\%，订单积压5140亿美元（环比+520亿，超预期7\%）。2026年资本支出上限上调150亿美元至2050亿美元，2027年预计3750亿美元。TPU销售路径清晰：2027年约3.5GW，2028年约4GW，毛利率约20\%。GenAI投入回报正在显性化。
{\small\color{refgray}数据：云业务Q2同比+82\%\par}
{\small\color{refgray}数据：EBIT利润率36\%\par}
{\small\color{refgray}数据：增量利润率54\%\par}
{\small\color{refgray}数据：订单积压5140亿美元\par}
{\small\color{refgray}边际变化：云业务高增长与高利润率同时出现，订单积压快速增长表明客户长期合同承诺加强，GenAI投入回报信号明确。\par}
\paragraph{MS} 中国自由流动性指标（M1增速-PPI-工业产出增速）6月降至2025年5月以来最低。M1仍低，PPI处高位，工业产出稳健。6月可能是PPI周期高点，下半年专项债发行加速将推动流动性改善，但企业借贷意愿偏弱制约改善幅度。该指标与MSCI中国指数同比变化存在相关性。
{\small\color{refgray}数据：自由流动性指标创2025年5月以来新低\par}
{\small\color{refgray}数据：M1增速仍处低位\par}
{\small\color{refgray}数据：PPI同比处高位\par}
{\small\color{refgray}数据：工业产出增速稳健\par}
{\small\color{refgray}边际变化：6月可能为PPI周期高点，下半年流动性有望改善但幅度有限，企业借贷意愿偏弱是核心制约。\par}
\paragraph{MS} 苹果六月季度iPhone收入预计超540亿美元（同比+22\%），受渠道补货和19\%出货增长驱动。九月季度展望分化：收入指引或略超预期，但毛利率指引（46.8\%）低于市场共识（47.4\%），服务收入增速低于共识130bps。当前估值约32x PE，短期积极因素已计入价格，财报需超预期才能推动上行。
{\small\color{refgray}数据：六月季度iPhone收入预计超540亿美元\par}
{\small\color{refgray}数据：iPhone出货量同比+19\%\par}
{\small\color{refgray}数据：九月季度毛利率估算46.8\% vs 市场共识47.4\%\par}
{\small\color{refgray}数据：服务收入增速低于共识130bps\par}
{\small\color{refgray}边际变化：市场已计入六月季度小幅超预期和九月季度收入上调，但毛利率和服务收入存在低于共识风险，短期交易空间有限。\par}
\paragraph{MS} 当前中国市场环境与2024年9月有本质区别：地方债务置换推进、核心CPI回升、房价调整放缓、上半年GDP增4.7\%在目标区间内。决策层不会推出类似2024年9月的大规模刺激，更可能在现有工具框架内微调。AI硬件出口贡献出口增量超10个百分点，亚洲工业资本开支超级周期提供第二支撑。
{\small\color{refgray}数据：上半年GDP增4.7\%\par}
{\small\color{refgray}数据：AI硬件出口贡献出口增量超10个百分点\par}
{\small\color{refgray}数据：亚洲除中国资本品进口增速创2004年以来新高\par}
{\small\color{refgray}数据：核心CPI从深度通缩回升至低通胀\par}
{\small\color{refgray}边际变化：与2024年9月多重风险自我强化不同，当前各维度均已触底或缓解，结构性出口为政策等待提供空间。\par}
\paragraph{MS} 欧央行7月按兵不动，但9月加息概率上升。触发条件：能源价格维持高位（天然气明显高于6月基线），或能源价格回落但PMI接近53且核心通胀上行。若能源价格显著下降且6月通胀预测过于悲观，9月才可能维持不变。市场对10月加息预期合理，但数据窗口更紧（仅两次PMI和一次HICP初值）。
{\small\color{refgray}数据：天然气价格明显高于欧央行6月基线\par}
{\small\color{refgray}数据：9月加息触发条件：能源高位或PMI接近53+核心通胀上行\par}
{\small\color{refgray}数据：10月会议前仅两次PMI和一次HICP初值可用\par}
{\small\color{refgray}边际变化：能源价格组合接近欧央行基线，9月加息成为基准情景而非尾部风险，除非能源价格显著回落。\par}
\paragraph{MS} 耐克收回线上分销授权对滔搏影响深远：线上渠道占F26收入22\%，且线上线下运营相互支撑。抖音渠道曾是线上销售占比从低双位数提升至35\%-40\%的主要驱动力。预计F27/F28收入连续下滑14\%，利润各降21\%，门店面积两年缩减约20\%，股息率或从135\%降至100\%。
{\small\color{refgray}数据：线上渠道占F26收入22\%\par}
{\small\color{refgray}数据：抖音渠道占线上销售35\%-40\%\par}
{\small\color{refgray}数据：F27/F28收入各降14\%\par}
{\small\color{refgray}数据：F27/F28利润各降21\%\par}
{\small\color{refgray}边际变化：耐克收回线上授权不仅削减收入池，还削弱线下门店流量来源和库存灵活性，可能形成循环压力。\par}
\subsubsection*{关键数据}
\begin{itemize}
  \item 30年期国债回报5.2\%，创2007年6月以来新高
  \item BofA牛熊指标9.6
  \item 中国对美出口占比下降7个百分点，对东盟上升5个百分点
  \item 西方脱钩成本约24万亿美元
  \item 日本产出缺口自2022年初持续为正
  \item 中国6月一般公共预算收入同比+8.7\%，土地出让收入-42.1\%
  \item AVC Q2毛利率预期32.5\% vs 市场一致预期约30\%
  \item 韩国造船P/B 0.5-0.6x，历史峰值约1.0x
  \item Alphabet云业务Q2同比+82\%，订单积压5140亿美元
  \item 中国自由流动性指标创2025年5月以来新低
\end{itemize}
\begin{figure}[htbp]
\centering
\includegraphics[width=0.88\linewidth]{figures/fig_048.jpg}
\caption{Exhibit 1 - Exhibit 1: MS China Free Liquidity Indicator vs. MSCI China YoY \% – Free liquidity declined to a fresh low since May 2025; we expect it to gradually improve into year-end MS ASIA LIMITED+ Laura Wang Equity Strategist}
\end{figure}
\begin{figure}[htbp]
\centering
\includegraphics[width=0.88\linewidth]{figures/fig_054.jpg}
\caption{Exhibit 3 - Exhibit 1: Tax revenue as a share of GDP appears to have stabilized after a prolonged decline Exhibit 2: Corporate and consumer sentiment indicators have largely stabilized above their troughs of late 2024 Exhibit 3: Core CPI h}
\end{figure}
\begin{figure}[htbp]
\centering
\includegraphics[width=0.88\linewidth]{figures/fig_059.jpg}
\caption{Exhibit 1 - Exhibit 1: We expect the ECB to hike rates in September ECB Deposit Rate: Mkt Pricing vs. MS Fcast (bp) Hold in July: As expected, the ECB remained on hold today. The policy paragraph remained largely unchanged (the reference to}
\end{figure}
\begin{figure}[htbp]
\centering
\includegraphics[width=0.88\linewidth]{figures/fig_036.jpg}
\caption{Figure 1 - Figure 1: Korea Shipbuilding: Share price performance \% Figure 2: Korea Shipbuilding: Foreign Ownership Changes Figure 3: Korea Shipbuilding: Consensus Revisions since 1Q26 earnings season - 2026E/2027E Sales}
\end{figure}
{\small\color{refgray}本节综合 12 篇研究；涉及 BofA、HSBC、JPM、MS。}
\newpage
\subsection{区域市场与主题：中国科技配置、能源供给与消费分化}
\textbf{全球机构资金正加速向中国科技板块迁移，同时能源市场面临多重供给扰动与结构性定价机会，消费领域则呈现区域与品类分化。}
\subsubsection*{主流共识}
\begin{itemize}
  \item 中国公募基金在2026年第二季度大幅增配科技板块，电子、通信持仓占比和超配比例均创历史新高，科技主题已成为核心布局方向。
  \item 全球原油市场短期受红海航运摩擦、哈萨克斯坦出口下降及夏季库存收紧支撑，价格上行风险增加。
  \item 雀巢二季度有机增长符合预期，但北美宠物护理库存调整拖累内生增长，AOA区域表现强劲而中国尚未贡献增量。
\end{itemize}
\subsubsection*{分歧与条件差异}
\begin{itemize}
  \item 中国科技板块的高持仓是否可持续：UBS认为资金迁移是系统性调整，但Bernstein对Spotify定价空间的乐观预期暗示科技消费仍有增长潜力，而NOM对中国运动品牌二季度增速放缓的观察则提示消费需求趋于谨慎。
  \item 美联储7月会议政策路径：NOM预计维持利率不变，但内部存在分歧，部分委员可能主张加息，市场对加息概率的预期可能过度。
  \item 能源价格中期走势：GS维持2027年布伦特原油均价75美元/桶的预测，但强调短期上行风险显著，而2027年可能面临供应过剩。
\end{itemize}
\subsubsection*{机构观点}
\paragraph{NOM} 中国运动品牌二季度销售增速普遍放缓，李宁与特步核心品牌同比转负，安踏主品牌和FILA增速收窄至低个位数。消费者对价值感要求提高，品牌间产品差异化不足，运营效率成为维持份额的关键。折扣力度加大但库存整体可控，特步渠道库存周转月数升至4.5-5个月。
{\small\color{refgray}数据：安踏主品牌二季度同比增速从一季度的中高个位数降至低个位数\par}
{\small\color{refgray}数据：李宁整体平台（不含童装）二季度同比低个位数下降\par}
{\small\color{refgray}数据：特步核心品牌二季度同比中个位数下降\par}
{\small\color{refgray}数据：安踏主品牌线下折扣约28\%，李宁线下折扣约35\%\par}
{\small\color{refgray}边际变化：二季度增速较一季度明显放缓，折扣力度加大但库存未显著恶化，显示品牌主动调整旧货而非被动积压。\par}
\paragraph{NOM} 美联储7月FOMC会议将维持政策利率3.625\%不变，会后声明无实质性调整。6月CPI低于预期，核心PCE月率可能降至0.2\%，为2025年11月以来首次。主席沃什新闻发布会将继续强调抗通胀信誉，但不会给出明确政策信号。达拉斯联储主席Logan和克利夫兰联储主席Hammack可能投反对票主张加息。
{\small\color{refgray}数据：6月核心PCE月率预计从5月的0.3\%降至0.2\%\par}
{\small\color{refgray}数据：BEA方法论调整将使核心PCE同比降低约20个基点\par}
{\small\color{refgray}数据：预计四季度核心PCE同比降至3.2\%（调整后约3.0\%）\par}
{\small\color{refgray}数据：市场预期约40\%概率上调利率\par}
{\small\color{refgray}边际变化：通胀数据边际改善叠加统计方法调整，为美联储维持观望提供数据支撑，但内部鹰派委员可能投反对票。\par}
\paragraph{UBS} 2026年第二季度中国主动管理型公募基金在大科技板块持仓占比达57.3\%，超配比例18.5\%，均创历史新高，远超此前大消费板块43.7\%的峰值。科创板与创业板持仓占比环比分别上升9.9和3.5个百分点，均达历史最高。跨境资金二季度净流入约2230亿元，扭转上季度净流出。
{\small\color{refgray}数据：大科技板块持仓占比57.3\%，超配比例18.5\%\par}
{\small\color{refgray}数据：电子板块持仓占比环比上升20.2个百分点\par}
{\small\color{refgray}数据：通信板块持仓占比环比上升4.4个百分点\par}
{\small\color{refgray}数据：科创板持仓占比环比上升9.9个百分点\par}
{\small\color{refgray}边际变化：科技主题从阶段性热点演变为公募核心布局，资金迁移是系统性调整而非边际调仓。\par}
\paragraph{Bernstein} Spotify定价空间尚未用尽，至少可延续至2027年。中期年化收入增速约15\%，其中用户增长贡献5\%，定价提升贡献6-7\%，增值服务贡献3-3.5\%。当前可寻址市场约1070亿美元，2030年将增至约1270亿美元。YouTube和Apple Music已跟进调价，全球价格与渗透率仍有提升空间。
{\small\color{refgray}数据：Spotify Premium当前可寻址市场约1070亿美元\par}
{\small\color{refgray}数据：2030年可寻址市场预计增至约1270亿美元\par}
{\small\color{refgray}数据：中期年化收入增速约15\%（定价贡献6-7\%）\par}
{\small\color{refgray}数据：全球付费音乐用户约8.6亿，Spotify占34\%份额\par}
{\small\color{refgray}边际变化：定价提升成为增长主驱动，增值服务产品（如HiFi、AI歌单）是弹性最大的变量，区域扩张以拉美和亚洲为主战场。\par}
\paragraph{GS} 自然资源公司2027年一致预期存在系统性修正空间。炼油板块MPC的2027年EPS预期高出市场7\%，受益于PADD 5区域结构性偏紧和出口能力扩大。中游LB的EBITDA预期高出市场7\%，核心驱动来自关联管道保底量协议爬坡。电力TLN的EBITDA预期高出市场7\%，受PJM电价走强和资产整合推动。
{\small\color{refgray}数据：MPC 2027年EPS预期高出市场约7\%\par}
{\small\color{refgray}数据：MPC 2026/2027年自由现金流回报约12\%/9\%\par}
{\small\color{refgray}数据：LB 2027年EBITDA预期高出市场7\%\par}
{\small\color{refgray}数据：LB每英亩收入预计从2026年约700美元增至2027年约921美元\par}
{\small\color{refgray}边际变化：市场对结构性因素（区域供应偏紧、管道爬坡、电价走强）的定价不足，分析师模型与共识存在系统性偏离。\par}
\paragraph{GS} 维持2026年四季度布伦特原油80美元/桶预测，但短期上行风险显著增加。红海航运摩擦、哈萨克斯坦出口下降及夏季库存收紧共同推升价格。曼德海峡每日近900万桶流量中约400万桶难以绕行。预计2027年布伦特均价75美元/桶，但全球战略储备采购需求（约120万桶/日）提供底部支撑。
{\small\color{refgray}数据：曼德海峡日均石油流量近900万桶，其中约400万桶难以绕行\par}
{\small\color{refgray}数据：沙特延布港原油装载量维持约500万桶/日\par}
{\small\color{refgray}数据：全球可见石油库存降至年内低点\par}
{\small\color{refgray}数据：2027年布伦特原油均价预测75美元/桶\par}
{\small\color{refgray}边际变化：短期价格上行风险显著增加，但2027年供应过剩预期不变，价格底部明确。\par}
\paragraph{GS} 雀巢二季度有机销售增长3.7\%，内生增长1.8\%，符合预期。北美RIG为-0.6\%，受宠物护理渠道库存调整拖累，但终端动销加速。欧洲RIG为0.5\%，受产品下架影响。AOA区域增长强劲，中国仍处调整期但降幅趋稳。管理层维持全年指引，中期17\%利润率目标未调整。
{\small\color{refgray}数据：二季度有机销售增长3.7\%，内生增长1.8\%\par}
{\small\color{refgray}数据：北美RIG为-0.6\%（宠物护理库存调整）\par}
{\small\color{refgray}数据：欧洲RIG为0.5\%（产品下架影响）\par}
{\small\color{refgray}数据：UTOP利润率16.4\%（含30bp一次性养老金收益）\par}
{\small\color{refgray}边际变化：北美宠物护理库存调整是短期扰动，终端动销已加速；下半年面临更难的RIG同比基数。\par}
\paragraph{GS} 三星Galaxy Z Fold/Flip 8系列定价上涨5-8\%（100-200美元），但DRAM与NAND配置与前代一致，涨价反映存储芯片成本传导而非规格升级。三星2026年智能手机出货量预计同比下降3\%至2.32亿部。MX/NW业务营业利润预计同比下滑93\%，但存储业务利润增长足以覆盖，推动整体营业利润同比增长8.6倍。
{\small\color{refgray}数据：Fold 8 Ultra起售价从1,999美元升至2,099美元\par}
{\small\color{refgray}数据：Flip 8起售价从1,099美元升至1,199美元\par}
{\small\color{refgray}数据：DRAM与NAND配置与前代一致\par}
{\small\color{refgray}数据：2026年智能手机出货量预计同比下降3\%至2.32亿部\par}
{\small\color{refgray}边际变化：利润重心从终端设备向存储芯片迁移，折叠屏提价是成本被动转嫁而非产品力驱动。\par}
\paragraph{GS} 安川电机5月ERP系统切换导致生产中断，各业务利用率从5月低点回升。AC伺服从51\%升至约70\%，逆变器从71\%升至接近100\%，机器人从34\%升至约70\%。公司维持季度营业利润160亿日元目标，认为一季度损失的利润可通过后续集中确认追回。AI相关需求强劲，伺服电机半导体设备订单同比增3倍。
{\small\color{refgray}数据：AC伺服利用率从5月51\%升至7月中旬约70\%\par}
{\small\color{refgray}数据：逆变器利用率从5月71\%升至7月接近100\%\par}
{\small\color{refgray}数据：机器人利用率从5月34\%升至7月约70\%\par}
{\small\color{refgray}数据：一季度营业利润85亿日元（内部目标120亿日元）\par}
{\small\color{refgray}边际变化：ERP事件影响有序恢复，AI需求对工业自动化设备订单拉动显著，利润损失本质是时间差而非订单取消。\par}
\subsubsection*{关键数据}
\begin{itemize}
  \item 中国公募基金大科技板块持仓占比57.3\%，超配比例18.5\%
  \item 曼德海峡日均石油流量近900万桶，约400万桶难以绕行
  \item Spotify当前可寻址市场约1070亿美元，2030年预计增至1270亿美元
  \item 雀巢二季度内生增长1.8\%，北美RIG为-0.6\%
  \item 三星Fold 8 Ultra起售价升至2,099美元，内存配置未变
  \item 安川电机机器人利用率从5月34\%回升至7月约70\%
  \item 美联储6月核心PCE月率预计降至0.2\%
  \item MPC 2027年EPS预期高出市场7\%
\end{itemize}
\begin{figure}[htbp]
\centering
\includegraphics[width=0.88\linewidth]{figures/fig_091.jpg}
\caption{Figure 1 - Figure 1: MFs' position change in Q226 (QoQ) Figure 2: Q226 MFs' \% of overweight/underweight by sector (benchmark against CSI300) Figure 3: MFs' sector allocation}
\end{figure}
\begin{figure}[htbp]
\centering
\includegraphics[width=0.88\linewidth]{figures/fig_096.jpg}
\caption{Exhibit 18 - EXHIBIT 2: We estimate Spotify Premium's current TAM to be \textbackslash{}\textasciitilde{}\textbackslash{}\$107B with Developed Europe, North America, and Asia as the biggest markets Estimated current TAM penetration by music DSPs (annual, \textbackslash{}\$B) ■ Other DSPs ■ Spotify ■ Unpe}
\end{figure}
\begin{figure}[htbp]
\centering
\includegraphics[width=0.88\linewidth]{figures/fig_140.jpg}
\caption{Exhibit 1 - Exhibit 2: Combined Oil Flows From Saudi Arabia and Sudan via Bab-al-Mandab Rose to 4mb/d in 2026Q2 (Amid Redirection Mostly via the East-West Pipeline and the Yanbu Port) \textbackslash{}*In a scenario with Hormuz, Bab-al-Mandab, and Suez ship}
\end{figure}
\begin{figure}[htbp]
\centering
\includegraphics[width=0.88\linewidth]{figures/fig_145.jpg}
\caption{Exhibit 1 - Exhibit 1: We forecast 3.3\% organic sales growth in FY26 Nestlé: Organic Sales Growth, split by driver, over time Exhibit 2: The US is 32\% of sales of which 45\% is Pet Care Nestlé: 2025 sales by country and region Exhibit 3: Fo}
\end{figure}
{\small\color{refgray}本节综合 9 篇研究；涉及 NOM、UBS、Bernstein、GS。}
\newpage
\subsection{AI基础设施投资挤压传统软件支出，IBM短期承压但需求延迟非消失}
\textbf{企业资本支出正从传统软件采购转向AI基础设施、存储和内存，导致IBM等传统软件厂商短期业绩承压；但推迟的交易并未取消，下半年业绩改善取决于客户支出正常化节奏及IBM执行能力。}
\subsubsection*{主流共识}
\begin{itemize}
  \item 企业资本支出正从传统软件采购转向AI基础设施、存储和内存，导致软件厂商短期收入增速放缓。
  \item 被推迟的大型企业交易属于需求延迟而非需求消失，部分已在后续季度重新完成。
  \item 软件业务增长依赖并购驱动，核心业务有机增长接近零。
  \item 经常性收入部分仍保持稳健增长，交易型收入受冲击较大。
\end{itemize}
\subsubsection*{分歧与条件差异}
\begin{itemize}
  \item 客户支出从AI基础设施/存储/内存采购中恢复的时间点和节奏存在不确定性，是下半年业绩改善的关键变量。
  \item IBM管理层预计全年收入增长4\%-5\%，但BARC认为执行能力和客户行为正常化节奏仍是不可控因素。
\end{itemize}
\subsubsection*{机构观点}
\paragraph{BARC} IBM 2Q26营收172亿美元，同比仅增1\%，低于市场预期约3\%，主因客户将资本支出从软件转向AI基础设施、存储和内存，导致大型企业交易推迟。软件收入77.6亿美元，同比增5\%，低于预期的80亿美元，增长由HashiCorp和Confluent收购驱动，核心业务有机增长约0\%。管理层将全年营收增长指引从5\%以上下调至4\%-5\%，软件增长指引从10\%以上下调至6\%-8\%。约三分之一推迟交易已在Q3前几周完成，但客户支出正常化节奏仍是下半年不可控变量。
{\small\color{refgray}数据：IBM 2Q26总营收172亿美元，同比增1\%，低于市场预期约3\%\par}
{\small\color{refgray}数据：软件收入77.6亿美元，同比增5\%，低于预期的80亿美元\par}
{\small\color{refgray}数据：核心软件业务有机增长约0\%，增长由HashiCorp和Confluent收购驱动\par}
{\small\color{refgray}数据：事务处理软件收入同比下滑9\%\par}
{\small\color{refgray}边际变化：客户支出从软件采购转向AI基础设施的结构性变化，导致IBM短期业绩低于预期并下调全年指引；但推迟交易部分恢复，显示需求延迟而非消失。\par}
\subsubsection*{关键数据}
\begin{itemize}
  \item IBM 2Q26总营收172亿美元，同比增1\%，低于市场预期约3\%
  \item 软件收入77.6亿美元，同比增5\%，低于预期的80亿美元
  \item 核心软件业务有机增长约0\%
  \item 事务处理软件收入同比下滑9\%
  \item 全年营收增长指引从5\%以上下调至4\%-5\%
  \item 软件增长指引从10\%以上下调至6\%-8\%
  \item 约三分之一推迟交易已在Q3前几周完成
  \item 约80\%软件收入为经常性收入，年化增长率约8\%
\end{itemize}
{\small\color{refgray}本节综合 1 篇研究；涉及 BARC。}
\newpage
\section{辅助信号｜战略咨询}
本板块整合波士顿咨询当日研究，聚焦企业韧性对长期回报的贡献、成本管理优先级上升及执行短板、欧洲太空探索的经济乘数效应、可再生能源扩张的核心驱动因素，以及AI在炼油行业的利润提升潜力。
\subsection{企业韧性：扰动时期的超额回报来源}
\textbf{BCG基于25年数据量化显示，扰动季度仅占研究周期的11\%，却解释了长期相对TSR的约30\%，韧性是长期超额回报的关键来源。}
\begin{itemize}
  \item 扰动季度内行业TSR分化程度（第25-75百分位差距30个百分点）是平稳期（16个百分点）的近两倍，扰动是区分优秀与普通公司的关键窗口。
  \item 韧性拆解为三个维度：降低即时冲击（贡献63\%）、恢复速度（23\%）、恢复程度（15\%），抗压能力比事后追赶更重要。
  \item 15\%的企业具备“通用韧性”（在超80\%扰动季度跑赢行业），这些公司25年年化TSR比行业高出5个百分点。
  \item 通用韧性企业的财务缓冲特征包括低负债，暗示财务储备与供应链冗余是扰动前缓冲设计的关键。
\end{itemize}
\begin{figure}[htbp]
\centering
\includegraphics[width=0.88\linewidth]{figures/fig_195.jpg}
\caption{Exhibit 5 - Outperformance is driven primarily by withstanding the immediate impact. Of the three phases of resilience, not all are equally important. Reducing the immediate impact of a shock accounts for 63\% of relative TSR in crisis periods, while greater recovery speed}
\end{figure}
\begin{figure}[htbp]
\centering
\includegraphics[width=0.88\linewidth]{figures/fig_197.jpg}
\caption{EXHIBIT 5 - Being generally resilient is extremely valuable over the long run, as such companies achieved an annual 5-pp TSR outperformance over their industries for the 25-year period that we observed (1995–2020). And, perhaps not surprisingly, five of the current Fortun}
\end{figure}
\subsection{成本管理优先级上升与执行短板}
\textbf{2025年成本管理取代增长成为企业第一战略优先级（33\%提及率，+8pp），但企业平均仅实现48\%的成本节约目标，文化阻力是最大障碍。}
\begin{itemize}
  \item 70\%的高管有足够中期可见度做出研究决策，表明企业是在有选择地收紧而非盲目调整。
  \item 未能实现长期结构性成本削减的公司中，最大障碍是“员工与组织文化”阻力，而非技术或资源不足。
  \item 主动推动文化对齐与敏捷管理的企业，成本削减效率比同行高出最多11个百分点。
  \item 供应链优化与产品组合简化是优先方向，67\%的高管计划将节省成本重新投入增长与创新。
\end{itemize}
\subsection{欧洲太空探索的经济乘数效应}
\textbf{BCG与欧洲空间政策研究所联合研究显示，2025-2040年500亿欧元太空探索计划可产生至少2600亿欧元累计GDP影响，乘数效应超5倍。}
\begin{itemize}
  \item 直接层面：新增约4000个太空岗位，产生250亿欧元直接GDP贡献；间接与诱导层面分别贡献700亿和550亿欧元，合计约1500亿欧元（3倍乘数）。
  \item 催化效应（新太空市场爆发）贡献额外1100亿欧元，是拉开乘数差距的核心来源。
  \item 欧洲目前仅占全球政府太空预算约15\%，与其约占全球25\%的GDP份额不匹配；全球太空经济2022年约4600亿美元，预计2040年达1万亿美元。
  \item 报告强调不参与的成本已不可忽视：欧洲若不行动，将损失全部宏观环境收益，且在发射、制造和在轨服务市场中份额可能持续调整。
\end{itemize}
\begin{figure}[htbp]
\centering
\includegraphics[width=0.88\linewidth]{figures/fig_199.jpg}
\caption{Figure 1 - Space exploration provides distinct levels of benefits that stem from the direct contribution of its core activities and propagate up to the level of the broader economy and society (Figure 1). Benefits of space exploration arise from a dedicated investment in}
\end{figure}
\begin{figure}[htbp]
\centering
\includegraphics[width=0.88\linewidth]{figures/fig_200.jpg}
\caption{Figure 2 - \#\# 3.1 Economic and societal benefits of space exploration Space exploration has the potential to create at least \textbackslash{}\textasciitilde{}€260 billion of cumulative GDP impact and an average of \textbackslash{}\textasciitilde{}90,000 FTEs in Europe between 2025 and 2040 (Figure 2). The investment in space explor}
\end{figure}
\subsection{可再生能源扩张的核心驱动：土地与现有装机基础}
\textbf{BCG对中美10年数据（超10万数据点）分析显示，廉价土地和现有装机基础是驱动风能/太阳能部署的首要因素，其重要性超过资源禀赋或国家目标。}
\begin{itemize}
  \item 廉价土地解释了中美约20-21\%的装机增长；美国土地成本最低的十个州装机增幅（24\%）是最高十个州（9\%）的2.6倍。
  \item 现有装机基础解释了约12-15\%的增长，体现经验曲线正向效应：早期部署降低后续供应链成本、施工经验和融资门槛。
  \item 美国需求增长解释了13\%的扩张，而中国增长动力更多来自国家能源转型战略和跨区输电规划，与本地需求关联微弱。
  \item 美国州级可再生能源目标在短期与装机增长相关，但长期无可测量影响，暗示政策目标本身不足以成为持续增长驱动力。
\end{itemize}
\begin{figure}[htbp]
\centering
\includegraphics[width=0.88\linewidth]{figures/fig_204.jpg}
\caption{Exhibit 1 - A close look at renewables at the province and state level reveals some surprises. \$\textasciicircum{}\{2\}\$ For example, even though China is the global leader in overall renewable energy production, the top five regions by share of electricity from renewables at the end of 202}
\end{figure}
\begin{figure}[htbp]
\centering
\includegraphics[width=0.88\linewidth]{figures/fig_205.jpg}
\caption{Exhibit 2 - \#\# EXHIBIT 2 The Top Drivers Influencing Wind and Solar Scaling in the US and China Are Similar Relative importance of wind and solar adoption drivers across US states and Chinese provinces Note: Drivers are listed in descending order by combined share of scal}
\end{figure}
\subsection{AI在炼油行业的利润提升潜力}
\textbf{BCG报告指出，炼油行业面临长期需求下降与产能过剩的结构性矛盾，AI到2030年可为中等水平玩家带来每桶0.5-1.2美元EBIT提升（超50\%利润改善）。}
\begin{itemize}
  \item 2026-2030年净炼油产能增长将超过需求增长，即使保守情景下供应仍超最乐观需求情景；航空煤油和石脑油需求更具韧性，汽油因电动化面临侵蚀。
  \item AI应用覆盖全价值链：规划与排程（每桶0.15-0.30美元）、运营（0.1-0.5美元）、物流与贸易（0.2-0.3美元）、维护与可靠性（0.05-0.1美元）。
  \item 案例显示，一家东南亚炼油企业通过AI驱动的数字孪生进行短期排程优化，实现每桶0.15-0.30美元收益。
  \item 多重因素（中国出口配额、全球增长与关税、俄罗斯制裁演变、新产能加速投产、IMO排放标准收紧）放大利润波动，AI成为从周期性孤岛优化转向实时利润决策闭环的核心杠杆。
\end{itemize}
\newpage
\section{辅助信号｜智库/国际机构}
本日智库/国际机构研究覆盖三大主题：亚太国防支出乘数、人民币计价结算条件、以及全球能源消费的“加法”逻辑。核心发现包括：国防支出在新兴市场产生温和持久正向产出效应（乘数0.7-1\%），但效果取决于财政空间与制度质量；人民币计价扩张依赖贸易依赖度与制度安排，而非对美元的直接替代；全球能源与材料消费遵循“加法”而非“替代”模式，清洁能源占比上升并未导致化石燃料绝对消费下降。
\subsection{亚太国防支出：温和乘数，条件依赖}
\textbf{国防支出在新兴市场（尤其亚太）产生温和但持久的正向产出效应，但效果高度依赖初始条件——财政空间、收入能力、制度质量与冲突暴露程度。}
\begin{itemize}
  \item 亚洲开发银行研究显示，1990-2023年间，新兴市场国防支出冲击在8年后产出乘数达0.7\%-1\%，而发达经济体效应偏弱甚至为负。
  \item 亚太地区国防支出增速显著高于其他新兴区域：2023年全球军事支出创纪录达2.4万亿美元，亚太新兴地区贡献超三分之一增量；中国支出2960亿美元，占区域近40\%。
  \item 越南、菲律宾过去十年国防预算增幅超50\%，印尼、马来西亚国防负担升至GDP的1.2\%-1.5\%，与拉美（1.3\%）和撒哈拉以南非洲（降至1.6\%）形成对比。
  \item 乘数差异源于初始条件：财政空间充裕、制度质量高、冲突暴露度低的地区，国防支出对产出的拉动更显著；反之则效果有限或为负。
  \item 研究采用局部投影方法识别经周期调整的国防支出冲击，确认核心结论在多种稳健性检验下保持稳定，但指出资源向基础设施、健康或教育领域配置可能带来更大长期收益。
\end{itemize}
\begin{figure}[htbp]
\centering
\includegraphics[width=0.88\linewidth]{figures/fig_164.jpg}
\caption{Figure 1 - \#\# 4.3. Stylized Facts Figure 1. Defense Spending as a Share of Gross Domestic Product: Median and Interquartile Range Across Country Groups (1990–2023) Figure 2. Composition of Government Defense Spending by Function (Share of Total Defense), 1990–2020}
\end{figure}
\subsection{人民币计价扩张：贸易依赖与制度安排驱动}
\textbf{人民币在国际贸易中计价结算的扩张并非对美元的直接替代，而是在美元主导框架下，通过贸易依赖、制度安排和政策工具共同推动的渐进过程。}
\begin{itemize}
  \item 亚洲开发银行研究所基于70个地区2016-2023年贸易计价数据发现，对中国的出口依赖度是解释人民币计价份额最稳定的变量，进口端效应比出口端更显著，反映中国作为买方的议价能力。
  \item 实行外汇管制、资本管制的经济体反而更倾向使用人民币而非美元，接入CIPS、签订货币互换协议显著提升人民币计价比例。
  \item 中国高层官员互访频率与政治制度联系强度也是显著变量，表明制度安排在推动人民币计价中发挥关键作用。
  \item 2016-2023年样本跨越中美贸易摩擦、新冠疫情及俄乌冲突后资金制裁常态化，外部冲击并未根本改变美元主导地位，但改变了部分地区在计价货币选择上的边际行为。
  \item 报告强调，贸易依赖度本身不充分——许多对华贸易依赖度高的地区人民币计价比例仍低，存在结构性障碍限制贸易关系向计价货币选择的传导。
\end{itemize}
\begin{figure}[htbp]
\centering
\includegraphics[width=0.88\linewidth]{figures/fig_169.jpg}
\caption{Figure 1 - Figure 1: As the local currency swap amount increases, so too does the renminbi settlement amount. However, since these settlement amounts include capital transaction figures, and as the data are only published for the top dozen o}
\end{figure}
\subsection{全球能源消费的“加法”逻辑：效率提升不必然减少用能}
\textbf{全球能源与材料流动历史上遵循“加法”而非“替代”逻辑，清洁能源占比上升并不必然导致化石燃料绝对消费量下降，效率提升反而可能刺激更多需求。}
\begin{itemize}
  \item IMF工作论文提出“广义杰文斯悖论”：自2012年以来，全球化石燃料消费总量与排放总量仍在增长，即便可再生能源在电力结构中份额持续扩大。
  \item 20世纪和21世纪的技术创新未导致任何主要能源或材料的绝对消费下降（唯一例外是羊毛、石棉和汞），所有能源载体和工业材料全球消费总量均呈上升趋势。
  \item 可再生能源基础设施建设本身需要大量钢铁、水泥和化工产品，而这些产品生产目前仍高度依赖化石燃料，形成“能源叠加”而非“替代”。
  \item 效率提升使能源变得更便宜、更容易获取，反而刺激了更多需求——全球数据清晰显示资源使用效率与总资源使用量之间呈强正相关。
  \item 报告将市场力量、地缘政治格局、贸易安排和资金因素纳入分析框架，解释为何全球能源与材料消费呈现结构性粘性，对资产叙事中“能源转型”的线性假设构成挑战。
\end{itemize}
\begin{figure}[htbp]
\centering
\includegraphics[width=0.88\linewidth]{figures/fig_189.jpg}
\caption{Figure 1 - \#\# Energy Transition? The data show an absence of energy transition in absolute terms. Variations in the shares of global sources of energy over time mask increasing consumption levels for all energies. While there have been large shifts in the shares of diffe}
\end{figure}
\subsection{宏观政策与金融稳定}
\textbf{用于校准利率、通胀、财政约束和金融稳定叙事。}
\begin{itemize}
  \item 爱尔兰作为小型开放地区，正同时面对贸易政策调整、人工智能扩散和能源转型三重结构性力量。IMF最新发布的国别报告（SIP/2026/070）通过可计算一般均衡模型，量化了这三股力量对爱尔兰产出、出口、就业和碳排放的长期影响。核心发现是：贸易政策变化主要带来贸易流向的重新配置，对总体...
  \item 中国在战略行业的补贴力度与贸易效应，已显著超过美国和欧盟。IMF最新工作论文通过构建多国多部门一般均衡模型，首次系统量化了2015至2023年间全球工业补贴对贸易流动和福利的影响。报告的核心发现是：中国对战略行业的补贴使其出口增长约12\%，而美国和欧盟的同类补贴仅带来约2\%的出口...
  \item 国际货币产品组织（IMF）2026年4月发布的技术援助报告，详细记录了为蒙古央行（BoM）开发并部署的现时预测与短期预测系统（NNTF）。该系统的核心价值在于：在验证测试中，其对关键指标的预测误差较历史人工预测降低了最高25\%，尤其在现时预测和未来一个季度的预测窗口上表现突出。这...
  \item 2025年11月和2026年2月，IMF通过高加索、中亚和蒙古区域能力发展中心（CCAMTAC）分两阶段对阿塞拜疆财政部开展了技术援助，核心目标是帮助该国建立按季度和年度编制并发布政府财政统计（GFS）与公共部门债务统计（PSDS）的能力。报告显示，阿塞拜疆已能定期向IMF提交年...
  \item 2015年11月16日，匈牙利福林（HUF）正式接入持续联系结算系统（CLS），成为中东欧地区首个通过支付对支付（PvP）机制完成外汇结算的货币。IMF三位研究员利用这一准自然实验，首次提供了结算不确定性在汇率定价中被定价的因果证据。
\end{itemize}
\begin{figure}[htbp]
\centering
\includegraphics[width=0.88\linewidth]{figures/fig_170.jpg}
\caption{Figure 1 - \#\# C. Tariff and Trade Agreements: Re-Shaping Global Trade 8. Trade policy shocks affect Ireland primarily through trade reallocation rather than through large aggregate output effects. Consistent with standard trade theory, the tariff shock has negative globa}
\end{figure}
\begin{figure}[htbp]
\centering
\includegraphics[width=0.88\linewidth]{figures/fig_175.jpg}
\caption{Figure 1 - Figure 1: Global exports of and subsidies to strategic products (2012-23 average) (b) Share of subsidy measures To evaluate the impact of industrial subsidies, we utilize a quantitative trade model incorporating external econo}
\end{figure}
\section{结语}
后续需验证的关键节点包括：中国财政政策下半年落地节奏及对企业借贷意愿的提振效果；欧洲央行9月加息概率及能源价格走势；AI基础设施投资向平台和应用层迁移的收入兑现速度；以及企业IT支出从AI硬件向传统软件恢复的时间点和节奏。
\newpage
\section{覆盖概览}
投行/券商 91 篇；战略咨询 5 篇；智库/国际机构 8 篇。\par
投行覆盖：GS 29篇 / MS 21篇 / JPM 11篇 / Bernstein 9篇 / NOM 7篇 / Citi 6篇 / HSBC 4篇 / BARC 2篇 / BofA 1篇 / UBS 1篇\par
\section{Disclaimer}
{\scriptsize\color{refgray}
This community is a paid learning and information-sharing community created for general educational purposes only. The membership fee is charged solely for access to community discussions, educational content, organization, and operational costs. It is not an advisory fee, management fee, performance fee, brokerage fee, or compensation for personalized financial, investment, legal, tax, or professional advice.\par
I participate and share information solely as an individual and for educational discussion among community members. I am not acting as a financial adviser, investment adviser, broker, dealer, portfolio manager, legal adviser, tax adviser, accountant, fiduciary, or any other licensed professional adviser. Nothing shared in this community constitutes, and should not be interpreted as, financial advice, investment advice, legal advice, tax advice, a recommendation, solicitation, offer, endorsement, instruction, or guarantee to buy, sell, hold, short, trade, or invest in any security, cryptocurrency, fund, derivative, strategy, or other financial product.\par
All content shared in this community, including messages, discussions, links, files, charts, examples, market commentary, watchlists, AI-generated or AI-polished summaries, and third-party materials, is general, impersonal, and educational in nature. It does not take into account any individual's financial situation, investment objectives, risk tolerance, investment horizon, liquidity needs, tax status, legal restrictions, or suitability requirements. No content should be relied upon as suitable for any specific person, account, portfolio, or financial decision.\par
Information shared in this community may be incomplete, inaccurate, outdated, speculative, AI-generated, AI-edited, or based on third-party internet sources that have not been independently verified. I make no representation or warranty as to the accuracy, completeness, timeliness, reliability, or suitability of any information shared.\par
Financial markets involve substantial risk, including the possible loss of principal. Past performance, historical data, hypothetical examples, simulated results, backtests, personal opinions, or educational examples do not guarantee future results. Each member is solely responsible for conducting their own research, exercising independent judgment, and making their own decisions. Before making any financial, investment, legal, or tax decision, members should consult qualified and licensed professionals.\par
Participation in this community, payment of any membership fee, or communication with me or other members does not create any adviser-client, fiduciary, professional, contractual, agency, partnership, employment, or other advisory relationship. By joining or participating in this community, you acknowledge and agree that you are solely responsible for your own decisions, actions, transactions, profits, losses, risks, and consequences, and that I shall not be liable for any direct, indirect, incidental, consequential, financial, legal, tax, or other loss or damage arising from or related to any content, discussion, information, or material shared in this community.\par
}
\newpage
\begin{center}
{\Large 更多详情报告kcdesk.com}\par\vspace{0.8cm}
\includegraphics[width=0.58\linewidth]{\detokenize{../../prompts/zsxq_img.jpg}}
\end{center}
\end{document}
